\documentclass[11pt]{amsart}
\usepackage[utf8]{inputenc}
\usepackage{color}%

\RequirePackage[numbers]{natbib}
\usepackage[dvips]{epsfig}
\usepackage{graphicx}
\usepackage{latexsym}
\usepackage{cancel}
\usepackage{ulem}
\usepackage[colorlinks=true,citecolor=blue,urlcolor=blue]{hyperref}

%\usepackage[notref, notcite]{showkeys}
\usepackage{amsmath,amsthm,amsfonts,amssymb,amscd,amsbsy,dsfont,hyperref}

\usepackage[margin=3cm]{geometry}
\setlength{\oddsidemargin}{.5cm} 
\setlength{\evensidemargin}{.5cm}
\setlength{\textwidth}{15cm} 
\setlength{\textheight}{20cm}
\setlength{\topmargin}{1cm}

\usepackage{aquabox}
\usepackage{stmaryrd}
\usepackage{graphicx}
\usepackage{float}
\usepackage{centernot}

\makeatletter
\def\l@section{\@tocline{1}{12pt plus2pt}{0pt}{}{\bfseries}}% <- added

\def\@tocline#1#2#3#4#5#6#7{\relax
    \ifnum #1>-1% <- added
  \ifnum #1>\c@tocdepth % then omit
  \else
    \par \addpenalty\@secpenalty%
    \begingroup \hyphenpenalty\@M
    \@ifempty{#4}{%
      \@tempdima\csname r@tocindent\number#1\endcsname\relax
    }{%
      \@tempdima#4\relax
    }%
    \parindent\z@ \leftskip#3\relax \advance\leftskip\@tempdima\relax
    \rightskip\@pnumwidth plus4em \parfillskip-\@pnumwidth
    #5\leavevmode\hskip-\@tempdima #6\nobreak\relax
    \hfil\hbox to\@pnumwidth{\@tocpagenum{#7}}\par
    \nobreak
    \endgroup
  \fi
\fi}
\makeatother
\setcounter{tocdepth}{2}  


\definecolor{vert}{rgb}{0,0.58,0}

\newcommand{\MD}{\color{vert}}


\usepackage{enumitem}
\newcommand{\subscript}[2]{$#1 _ #2$}
\newcommand{\delete}[1]{}

\newcommand{\E}{\espe}
\newcommand{\R}{\RR}
\renewcommand{\Re}{\mathfrak{R}e}

\newcommand{\norm}[1]{\lVert#1\rVert_2}

\newcommand{\rj}{\ln^2(J)}
\newcommand{\redline}{\noindent\textcolor{red}{\rule{16cm}{1mm}}}
\newcommand{\ul}[1]{\underline{#1}}
\newcommand{\blue}[1]{\textcolor{blue}{#1}}

\newcommand{\change}[1]{\textcolor{red}{#1}}

\newcommand{\mmax}{m_{\max}}
\newcommand{\tdiff}{(t_2 - t_1)}
\newcommand{\mopt}{m_J^\star}

\newcommand{\cte}{}

\newcommand{\diff}{\,\mathrm{d}}

\newenvironment{details}
  {\cor{violet}}%
  {}
\newenvironment{review}%
  {\color{red}}%
  {}
 \newenvironment{review1}%
  {\color{blue}}%
  {}

\newcommand{\ndr}[1]{\textcolor{blue}{[#1]}}

\newcommand{\supp}{\mathrm{supp}\,}
\newcommand{\spec}{\mathrm{spec}\,}
\newcommand{\tendsto}[1]{\xrightarrow[#1]{}}

\title{A mutation equation}
\author{}
\date{\today}

\begin{document}

\maketitle
\begin{center}
    \textbf{Version of the \today}
\end{center}

\tableofcontents

% \section{The Mutation Equation : all mutation are deleterious}

% \subsection{Description of the model}

% Consider the \textit{mutation equation}
% \begin{equation}\left\{
%     \begin{array}{ll}
% 	\frac{\partial}{\partial t}u(t,x) + \lambda u(t,x) = \lambda \int_{0}^{\infty} k(y, x) u(t,y) \dy
% 	\eqsepv
% 	& t \in (0, \infty), x \in \RR_+, \lambda \in (0, \infty) \eqfinv
% 	\\
% 	u(0, x) = u_0(x)
% 	& x \in \RR_+
% 	\eqfinp
%     \end{array}
% 	\right.
% 	\label{eq:pseudo}
% \end{equation}

% In all that follows, we will assume that $k$ is self-similar, \ie 
%     \begin{equation*}
%         k(y, x) = \frac{1}{y} k_0 \Big(\frac{x}{y} \Big)
%         \eqsepv
%         y > 0
%         \eqsepv 
%         x \in \RR_+
%         \eqfinp
%     \end{equation*}
%     \begin{equation*}
%         \supp k_0 \subset [0, 1] 
%         \eqsepv
%         \boxed{\int_0^1 k_0(z) \dz = 1}
%         \eqfinp
%     \end{equation*}

% In addition, we assume that 
%     \begin{equation}
%         \boxed{\int_0^1 k_0(z) (1 + z) \dz < \infty}
%         \eqfinp
%         \label{H:douesco4}
%     \end{equation}
% \onlyME{En cas de besoin, on supposera l'existence de moments pour tout $n \in \NN$.}
    
% \begin{remark}
%     This last assumption ensures that the Mellin transform $K(s) := \Mcal_{k_0}(s) = \int_0^\infty k_0(z)z^{s-1}\dz$ is well defined for all $\Re(s) \in (1, 2)$.
% \end{remark}
    
% Under theses assumptions, Equation\eqref{eq:pseudo} becomes
%     \begin{equation}
%     	\frac{\partial}{\partial t}u(t,x) + \lambda u(t,x) = \lambda \int_{x}^{\infty} \frac{1}{y} k_0 \Big(\frac{x}{y} \Big) u(t,y) \dy
% 	    \eqfinp 
% 	    \label{eq:similar_pseudo_growth_equation}
%     \end{equation}

% \subsubsection{Explicit solution to the mutation Equation \eqref{eq:similar_pseudo_growth_equation}}

% We assume that     \begin{equation}
%         u_0 \geq 0 \mtext{and} \int_0^\infty u_0(x) (1+x) \dx < \infty
%         \eqfinp
%         \label{H:u0_L1}
%     \end{equation}
    
% Define 
%     \begin{equation}
%         I(u_0) := \left\{ p \in \RR \eqsepv \int_0^\infty u_0(x) x^{p-1} \dx < \infty \right\}
%         \eqfinp
%     \end{equation}
% By Assumption~\eqref{H:u0_L1}, $[1, 2] \subset I(u_0)$. Moreover $I(u_0)$ is a interval of $\RR$.\\
% It leads us to define
%     \begin{subequations}
%     \begin{align*}
%         p_0 &= \inf \{p : p \in I(u_0) \} \in [-\infty, 1] \eqfinv\\
%         q_0 &= \sup \{q : q \in I(u_0) \} \in [2, \infty] \eqfinp
%     \end{align*}    
%     \end{subequations}

% For the simplicity, we assume that 
%     $ p_0=-\infty$ and $q_0 = \infty$
    
% We also assume that 
%     \begin{equation}
%         \forall \nu \in (p_0, q_0) 
%         \eqsepv
%         \int_0^\infty |u_0'(x)| x^\nu \dx < \infty
%         \eqsepv
%         \lim_{x \to 0} x^\nu u_0(x) = 0
%         \eqsepv
%         \lim_{x \to \infty} x^\nu u_0(x) = 0
%         \label{H:DoumicEsco22}
%     \end{equation}
%     \begin{equation}
%         \forall \nu \in (p_0, q_0) 
%         \eqsepv
%         \int_0^\infty |u_0''(x)| x^\nu \dx < \infty
%         \eqsepv
%         \lim_{x \to 0} x^{\nu + 1} u_0'(x) = 0
%         \eqsepv
%         \lim_{x \to \infty} x^{\nu + 1} u_0'(x) = 0
%         \label{H:DoumicEsco23}
%     \end{equation}
    
% In the following, we consider the Mellin transforms
%     \begin{equation*}
%         U(t,s) := \Mcal_{u(t, \cdot)}(s)
%         \mtext{and}
%         K(t,s) := \Mcal_{k_0(t, \cdot)}(s)
%     \end{equation*}
% when these integrals are well-defined.

% \begin{theorem} Let 
% $u(t, x) \in \Ccal([0, \infty); L^1((1+x)\dx)) \cap \Ccal^1([0, \infty); L^1((1+x)\dx))$ be a solution of Equation\eqref{eq:similar_pseudo_growth_equation}. Then 
% 	\begin{equation*}
% 		\frac{\partial}{\partial t} U(t,s) + \lambda U(t,s) = \lambda K(s) U(t,s)
% 		\eqfinp
% 	\end{equation*}
% \end{theorem}
% \begin{proof}
% Let multiply Equation\eqref{eq:similar_pseudo_growth_equation} by $x^{s-1}$ and integrate over $(0, \infty)$, we obtain
% 	    \begin{align*}
%     	   \frac{d}{dt} \int_0^\infty u(t,x) x^{s-1} \dx
%     	   + 
%     	   \lambda \int_0^\infty u(t,x)x^{s-1} \dx
%     	   &= 
%     	   \lambda \int_0^\infty x^{s-1} \int_{x}^{\infty} \frac{1}{y} k_0 \Big(\frac{x}{y} \Big) u(t,y) \dy \dx \\
% 	       &= 
% 	       \lambda \int_0^\infty  u(t,y) \int_{0}^{y} \frac{1}{y} k_0 \Big(\frac{x}{y} \Big) x^{s-1} \dx \dy 	   \\
% 	       &= 
% 	       \lambda \int_0^\infty  u(t,y) y^{s-1}  \int_{0}^{1} k_0(z) z^{s-1} \dz \dy	   \\
% 	       &= \lambda \int_0^\infty  u(t,y) y^{s-1} \dy  \int_{0}^{1} k_0(z) z^{s-1} \dz
% 	       \eqfinp
% 	    \end{align*}
% Then 
% 	\begin{equation*}
% 		\frac{\partial}{\partial t} U(t,s) + \lambda U(t,s) = \lambda K(s) U(t,s)
% 		\eqfinp
% 	\end{equation*}
	
% Then 
%     \begin{equation*}
%         U(t,s) = U_0(s) e^{\lambda(K(s) - 1) t}
%         \eqsepv
%         U_0(s) = \Mcal_{u_0}(s)
%         \eqfinp
%     \end{equation*}
% \end{proof}

% By using the inverse Mellin transform, we have at least formally the explicit solution

% \begin{equation}
%     u(t,x) = \frac{1}{2i\pi} \int_{c - i \infty}^{c + i \infty} U_0(s) e^{\lambda(K(s) - 1) t} x^{-s}\ds 
%     \eqsepv x \in (0, \infty) \as
%     \label{eq:explicit_reconstruction}
% \end{equation}

% \begin{theorem}(Existence of an explicit solution) Assume \eqref{H:u0_L1}, \eqref{H:DoumicEsco22} and \eqref{H:DoumicEsco23}.
% Then, there exists a unique function $u \in \Ccal([0, \infty); L^1((1+x)dx)) \cap \Ccal^1((0, \infty); L^1((1+x)dx))$ that satisfies Equation~\eqref{eq:similar_pseudo_growth_equation} for all $t > 0$ and almost every $x > 0$. This solution is given by \eqref{eq:explicit_reconstruction} for all $\nu \in (p_0, q_0)$.
% \end{theorem}
% \begin{proof}
% The result is a consequence of the Mellin inversion formula.
% In order to apply it, we have to prove that for any $\nu \in (1, 2)$, 
%     \begin{equation*}
%         \int\limits_{\nu - i \infty}^{\nu + i \infty} \bigg| U_0(s) e^{\lambda(K(s) - 1) t} x^{-s} \bigg| \ds < \infty
%         \eqfinp
%     \end{equation*}
% We know that for all $s \in [1, 2]$, 
%     \begin{align*}
%         K(s) = \int_{0}^{\infty} k_0(z) z^{s-1} \dz
%         &= \int_{0}^{1} k_0(z) z^{s-1} \dz \\
%         &\leq \int_{0}^{1} k_0(z) \dz
%         \\
%         &= K(1)  \\
%         &= 1
%         \eqfinp
%     \end{align*}
    
% Then 
%     \begin{equation*}
%          \bigg| U_0(s) e^{\lambda(K(s) - 1) t} x^{-s} \bigg|  <   \bigg| U_0(s)  x^{-\nu} \bigg|
%     \end{equation*}

% It only remains to prove that 
%     \begin{equation}
%             \int\limits_{\nu - i \infty}^{\nu + i \infty} \big| U_0(s)\big| \ds 
%             = 
%             \int\limits_{-\infty}^{\infty} \big| U_0(\nu + is)\big| \ds
%             < \infty
%             \label{eq:DouEsc43}
%     \end{equation}

% This quantity has already been studied under the same assumptions by Doumic and Escobedo \cite{doumic2016time}. Integration by parts lead to 
%     \begin{align*}
%         U_0(\nu + i v) 
%         &= \int_0^\infty u_0(x) x^{\nu-1}x^{i v}\dx \\
%         &= \frac{1}{iv + 1} \big[ u_0(x) x^{\nu} x^{i v} \big]_0^\infty - \frac{1}{iv + 1} \int_0^\infty (u_0(x) x^{\nu-1})' x^{i v+1}\dx \\
%         \eqfinp
%     \end{align*}
% Assumption \eqref{H:DoumicEsco22} guarantees that $[ u_0(x) x^{\nu} x^{i v} ]_0^\infty = 0$. Therefore
%     \begin{align*}
%         U_0(\nu + i v) 
%         &= - \frac{1}{iv + 1} \int_0^\infty (u_0(x) x^{\nu-1})' x^{i v+1}\dx \\
%         &= - \frac{\nu - 1}{iv + 1} \int_0^\infty u_0(x) x^{\nu-2} x^{i v + 1}\dx  - \frac{1}{iv + 1} \int_0^\infty u_0'(x) x^{\nu-1} x^{i v + 1}\dx \\
%         &= - \frac{\nu - 1}{iv + 1} \int_0^\infty u_0(x) x^{\nu-1} x^{i v}\dx  - \frac{1}{iv + 1} \int_0^\infty u_0'(x) x^{\nu} x^{i v}\dx \\
%         &= - \frac{\nu - 1}{iv + 1} U_0(\nu + i v)   - \frac{1}{iv + 1} \int_0^\infty u_0'(x) x^{\nu} x^{i v}\dx
%         \eqfinp
%     \end{align*}
% It follows that 
%     \begin{align*}
%         U_0(\nu + i v)
%         &= - \frac{1}{(iv + \nu)}\int_0^\infty u_0'(x) x^{\nu} x^{i v}\dx
%         \eqfinp
%     \end{align*}
% \onlyME{It follow that 
%     \begin{equation*}
%         |U_0(\nu + i v)| \leq \frac{1}{\sqrt{\nu^2 + v^2}} \int_0^\infty |u_0'(x)| x^{\nu} \dx
%         \eqfinp
%     \end{equation*}}
% The same calculations and Assumption\eqref{H:DoumicEsco23} allow us to obtain 
%     \begin{align*}
%         U_0(\nu + i v)
%         &=  \frac{1}{(iv + \nu)(iv + \nu + 1)}\int_0^\infty u_0''(x) x^{\nu + 1} x^{i v}\dx
%         \eqfinp
%     \end{align*}
% It follow that 
%     \begin{equation*}
%         |U_0(\nu + i v)| \leq \frac{1}{\sqrt{(\nu^2 + v^2)((1+\nu)^2 + v^2)}} \int_0^\infty |u_0''(x)| x^{\nu + 1} \dx
%         \eqfinp
%     \end{equation*}

% We have that $\int_0^\infty |u_0''(x)| x^{\nu + 1} \dx < \infty$ from Assumption\eqref{H:DoumicEsco23}. It allows us to deduce \eqref{eq:DouEsc43} \\

% Conversely, if $u_0$ satisfies \eqref{H:u0_L1}, \eqref{H:DoumicEsco22} and \eqref{H:DoumicEsco23}, then 
%  $$u \in \Ccal([0, \infty); L^1((1+x)dx)) \cap \Ccal^1((0, \infty); L^1((1+x)dx))$$
% where $u$ is defined by \eqref{eq:explicit_reconstruction}. By construction, it is a solution of the equation \eqref{eq:similar_pseudo_growth_equation}.
% \end{proof}

% \subsection{Speed of convergence}

% In the following, we define
%     \begin{equation*}
%         w(t,x) := e^t u(t,x)
%     \end{equation*}
% Then $\forall c \in (1, 2)$
%     \begin{equation}
%         w(t,x) = \frac{1}{2i\pi} \int_{c - i \infty}^{c + i \infty} U_0(s) e^{\lambda K(s) t} x^{-s}\ds 
%     \eqsepv x \in (0, \infty) \as
%     \end{equation}

% We define $s_+(t,x) := (K')^{-1} \Big( \frac{\log x}{\lambda t}\Big)$

% \begin{theorem}
%     Assume that $u_0$ satisfies the condition \ref{H:u0_L1}, \ref{H:DoumicEsco22} and \ref{H:DoumicEsco23} and let $\varepsilon(t)$ such that $\varepsilon(t)\to 0$ but $t\varepsilon^2(t) \to \infty$ as $t\to \infty$\\

%  Suppose that  the measure $k_0$ has a non zero absolutely continuous part. Then for all $\delta >0$ arbitrarily small, there exists two functions $\gamma _{ \delta  } (\varepsilon )$ and $\omega _{ \delta  } (\varepsilon )$ such that:

% \begin{eqnarray}
% &&u(t, x)=x^{-s_+(t, x)}e^{(K(s_+(t, x))-1)t}\times \nonumber\\
% &&\hskip 3cm \left(  \frac {U_0(s_+(t, x))+\omega _{ \delta  } (\varepsilon )} {\sqrt{2\pi t K''(s_+)}}\,\Theta_1(t)+\mathcal O\left(e^{- t\gamma  _{ \delta  }(\varepsilon (t)) }\right)  \right) \label{theorem3:1}\\
% &&\hskip 1.5cm \hbox{as}\,\,t\to \infty,\,\,\hbox{uniformly for}\; x\;\hbox{such that}\,\,  p_0+\delta<s _+(t, x)<q_0-\delta, \nonumber\\
% &&\Theta_1(t)=\left(1+
%  \mathcal O\left(\frac {e^{-\frac {t\varepsilon^2 (t)} {2}K''(q_0)}} {\sqrt {t \varepsilon^2 (t)K''(q_0)}}\right)\right),\,\,\,\hbox{as}\,\,t\to \infty\label{theorem3:2}
% \end{eqnarray}
% where,  $\lim _{ \varepsilon \to 0}\varepsilon  ^{-2}\gamma _{ \delta  }  (\varepsilon  )=C _{ \delta  }$ for some constant $C _{ \delta  }\in (K^{''} (q_0) /2 ,  K^{''} (p_0)/2)$. The function  $\omega _{ \delta  }(\varepsilon )$, defined later, satisfies
% \begin{eqnarray}
% \label{theorem3:omega}
% &&|\omega _{ \delta  } (\varepsilon )|\le C (p_0+\delta,q_0-\delta)\left( \int _{R(\varepsilon )}^\infty \left|z^{q_0-\delta-1}u_0(z)\right|dz+\right. \nonumber \\
% &&\hskip 4cm \left. +\int _0^{\frac {1} {R(\varepsilon )}} \left|z^{p_0+\delta-1}u_0(z)\right|dz+\varepsilon  \log R(\varepsilon )\right),
% \end{eqnarray}
% where the constant $C(p_0+\delta,q_0+\delta)$ depends on the integrability properties of $u_0$, $R$ is any function such that $R(\varepsilon )\to \infty$, $\varepsilon \log R(\varepsilon )\to 0$ as $\varepsilon \to 0$.
% \end{theorem}

% \subsubsection{Proof}

% Here, we want to apply the \textit{steepest descent method} \cite{debye1909naherungsformeln} when $t \to \infty$.

% We recall that 
%     \begin{equation*}
%         \omega(t,x)
%         = \int_{c-i\infty}^{c+i\infty} U_0(s) e^{\lambda K(s)t}x^{-s} \ds
%         =  \int_{c-i\infty}^{c+i\infty} U_0(s) e^{\phi(s, t, x)} \ds
%         \eqfinp
%     \end{equation*}
% with $\phi(s,t,x) := \lambda K(s) t - s \ln(x) $.\\

% In order to apply the steepest descent method, our first step consist in finding the stationary point of $\phi'(\cdot, t, x) = 0$. 

% Define $p_0 < s_+(t,x) := (K')^{-1} \Big( \frac{\log x}{\lambda t} \Big) < q_0$.

% \paragraph{\ul{Step 1 -- Show that $s \mapsto \Re(\phi(s, t, x))$ reachs its maximum at $s_+(t, x)$}} \text{}\\

% $\blacktriangleright$ Consider a neighborhood of $s_+(t,x)$ : $|s - s_+(t,x)| < \varepsilon$. We have
% \begin{equation}
%     K(s) = K(s_+) + (s - s_+(t,x)) K'(s_+(t,x)) + \frac{(s - s_+(t,x))^2}{2} K''(s_+(t,x)) + \frac{(s - s_+(t,x))^3}{6} K'''(c(t,x))
% \end{equation}
% where $s < c(t,x) < s_+(t,x)$.
%     \begin{align*}
%          |K'''(c)|
%          &= |\int_0^1 (\log z)^3 k_0(z) e^{(c-1)\log(z)}\dz|
%          \leq \int_0^1 |(\log z)^3| k_0(z)  e^{\Re (c-1)\log(z)}\dz \\
%          &\leq  \boxed{\int_0^1 |(\log z)^3| k_0(z)  e^{(q_0-1)\log(z)}\dz < \infty}
%     \end{align*}
    
% Given that $|s - s_+| < \varepsilon$ and $s_+ + \varepsilon < q_0$, we have 
%     \begin{equation}
%         \Re K(s) = K(s_+) + \frac{(s - s_+(t,x))^2}{2} K''(s_+(t,x)) + O(\varepsilon^3)
%         \eqfinp
%     \end{equation}
% and 
%     \begin{equation}
%         \Re \phi(s, t, x) = \phi(s_+, t, x) + \lambda t \Big(\frac{(s - s_+(t,x))^2}{2} K''(s_+(t,x)) + O(\varepsilon^3)\Big)
%         \eqfinp
%     \end{equation}

    
% $\blacktriangleright$ Outside the neighborhood of $s_+(t,x)$ : $|s - s_+(t,x)| < \varepsilon$. We have
%     \begin{align*}
%         \Re(\phi(s_+ + i \nu, t, x) - \phi(s_+, t, x)) 
%         &=
%         \Re \bigg( \int_0^1 x^{s_+ - 1}k_0(x) (x^{i\nu} - 1) \dx \bigg) \\
%         &=
%         \Re \bigg( \int_0^1 x^{s_+ - 1}k_0(x) (e^{i\nu \log(x)}- 1) \dx \bigg) \\ 
%         &= 
%         \int_0^1 x^{s_+ - 1}k_0(x) [\cos(\nu \log(x))- 1] \dx  \\
%         &=: G(v) \leq 0
%         \eqfinp
%     \end{align*}
    
% Recall the that $k_0(x) = g(x)dx + \dmu + \dnu$ where
% $g \in L^1(\dx)$
% $\dmu$ is a singular continuous measure and $\dnu$ a singular discrete measure.\\ 

% In the following, we assume that $k_0(x) :=  g(x)dx$.\\

% For any $\nu \in \RR \setminus \{0\}$, we have
% \begin{equation}
%      \int_0^1 x^{s_+ - 1}g(x) [\cos(\nu \log(x))] \dx < \int_0^1 x^{s_+ - 1}g(x) \dx
%      \eqfinp
% \end{equation}

% It follows that $G(\nu) < 0$, then there exists $\varepsilon > 0$, $R >0$ and $\delta > 0$
%  such that $\sup_{\varepsilon < \nu < R}G(\nu) < - \delta$.\\
 
%  It remains to verify that $\lim_{\nu \to +\infty}G(\nu) \neq 0$. 
%     \begin{align*}
%         G(\nu) 
%         &=  \int_0^1 x^{s_+ - 1}k_0(x) [\cos(\nu \log(x))- 1] \dx \\
%         &=  \int_{0}^\infty e^{-s_+ y}k_0(e^{-y})(\cos(-\nu y) - 1 )\dy
%         \end{align*}
% From Riemann–Lebesgue lemma, we have that 
%     \begin{equation}
%         \lim_{\nu \to +\infty}G(\nu) = - \int_{0}^\infty e^{-s_+ y}k_0(e^{-y})\dy = - \int_{0}^1 x^{s_+ - 1}k_0(x)\dx < 0
%     \end{equation}
    
% Therefore
% \begin{lemma} For all $\varepsilon > 0$, there exist $\gamma(\varepsilon)$ such
%     \begin{equation}
%         \sup_{|\nu| \geq \varepsilon } \Re(\phi(s_+ + i \nu, t, x)) \leq \Re(\phi(s_+, t, x)) - \gamma(\varepsilon)
%         \eqfinp
%     \end{equation}
% and
%     \begin{equation}
%         \sup_{|\nu| \leq \varepsilon } \Re(\phi(s_+ + i \nu, t, x)) \leq \Re(\phi(s_+, t, x))
%         \eqfinp
%     \end{equation}
% \end{lemma}

% \paragraph{\ul{Step 2 -- Use of the right path}} \text{}\\

% Close to $s_+(t,x)$, we consider the quantity
%     \begin{align*}
%     \frac{1}{2i \pi} \int_{s_+(t,x) - i\varepsilon}^{s_+(t,x) + i\varepsilon} U_0(s) e^{\phi(s,t,x)} \ds 
%     &=
%     \frac{1}{2i \pi} U_0(s_+(t,x)) \int_{s_+(t,x) - i\varepsilon}^{s_+(t,x) + i\varepsilon} e^{\phi(s,t,x)} \ds
%     \\
%     & \qquad \qquad \qquad+ 
%     \frac{1}{2i \pi} \int_{s_+(t,x) - i\varepsilon}^{s_+(t,x) + i\varepsilon} (U_0(s) -  U_0(s_+(t,x)))  e^{\phi(s,t,x)} \ds \\
%     &= I_1 + I_2
%     \end{align*}
    
% $\blacktriangleright$ \ul{The terme $I_1$}
% For any $s \in (s_+ - i \infty, s_+ + i \infty)$ such that $|s - s_+| < \varepsilon$
% We know that 
%     \begin{equation}
%         \Re \phi(s, t, x) = \phi(s_+, t, x) + t \Big(\frac{(s - s_+(t,x))^2}{2} K''(s_+(t,x)) + O(\varepsilon^3)\Big)
%         \eqfinp
%     \end{equation}

% Then 
% \begin{align*}
%    \int_{s_+(t,x) - i\varepsilon}^{s_+(t,x) + i\varepsilon} e^{\Re \phi(s,t,x)} \ds
%     &=
%     e^{\phi(s_+, t, x)} \cdot \int_{s_+(t,x) - i\varepsilon}^{s_+(t,x) + i\varepsilon} e^{t \Big(\frac{(s - s_+(t,x))^2}{2} K''(s_+(t,x)) + O(\varepsilon^3)\Big)} \ds \\
%     &= 
%      i e^{\phi(s_+, t, x)} \cdot \int_{- \varepsilon}^{\varepsilon} e^{t \Big(\frac{(s - s_+(t,x))^2}{2} K''(s_+(t,x)) + O(\varepsilon^3)\Big)} \ds \\
%     &= 
%      i e^{\phi(s_+, t, x)} \cdot \int_{- \varepsilon}^{\varepsilon} e^{- t \Big(\frac{v^2}{2} K''(s_+(t,x)) + O(\varepsilon^3)\Big)} \ds 
% \end{align*}

% Now we can see that we are back to the classical method of Laplace. The main idea is to transform the previous integral into a Gauss integral to have
%     \begin{equation}
%     \int_{- \varepsilon}^{\varepsilon} e^{- t \Big(\frac{v^2}{2} K''(s_+(t,x)) + O(\varepsilon^3)\Big)} \ds 
%     \sim_{t \to \infty}
%     \sqrt{2\pi} \frac{1}{\sqrt{tK''(s_+)}}
%     \end{equation}
    
% Indeed, 
% \begin{align*}
%     \int_{- \varepsilon}^{\varepsilon} e^{- t \Big(\frac{v^2}{2} K''(s_+(t,x)) + O(\varepsilon^3)\Big)} 
%     &= 
%      \frac{1}{\sqrt{tK''(s_+)}} \int_{- \varepsilon \sqrt{tK''(s_+) + \Ocal(\varepsilon})}^{\varepsilon \sqrt{tK''(s_+) + \Ocal(\varepsilon})} e^{- \frac{y^2}{2}}\dy \\
%      \intertext{and}
%      \int_{- \varepsilon \sqrt{tK''(s_+) + \Ocal(\varepsilon})}^{\varepsilon \sqrt{tK''(s_+) + \Ocal(\varepsilon})} e^{- \frac{y^2}{2}}\dy &= 
%     \sqrt{2\pi} - 2 e^{-\frac{t(\varepsilon^2 K''(s_+) + \Ocal(\varepsilon^3))}{2}} \cdot \Ocal\bigg(\frac{1 }{\varepsilon \sqrt{tK''(s_+) + \Ocal(\varepsilon})}\bigg)
%     \eqfinp
% \end{align*}
% It follows that 
%     \begin{align*}
%      \frac{1}{2i\pi}   \int_{s_+(t,x) - i\varepsilon}^{s_+(t,x) + i\varepsilon} e^{\Re \phi(s,t,x)} \ds 
%      &= 
%      \frac{1}{2i\pi} i e^{\phi(s_+, t, x)} \cdot \int_{- \varepsilon}^{\varepsilon} e^{- t \Big(\frac{v^2}{2} K''(s_+(t,x)) + O(\varepsilon^3)\Big)} \ds  \\
%      &= 
%       \frac{1}{2\pi} e^{\phi(s_+, t, x)} \cdot \frac{1}{\sqrt{tK''(s_+)}} \int_{- \varepsilon \sqrt{tK''(s_+) + \Ocal(\varepsilon})}^{\varepsilon \sqrt{tK''(s_+) + \Ocal(\varepsilon})} e^{- \frac{y^2}{2}}\dy \\
%      &=
%     \frac{1}{2\pi} e^{\phi(s_+, t, x)} \cdot \frac{1}{\sqrt{tK''(s_+)}} \bigg( 
%     \sqrt{2\pi} - 2 e^{-\frac{t(\varepsilon^2 K''(s_+) + \Ocal(\varepsilon^3))}{2}} \cdot \Ocal\bigg(\frac{1 }{\varepsilon \sqrt{tK''(s_+) + \Ocal(\varepsilon})}\bigg) \bigg) \\
%     \end{align*}
    
% Therefore, 
%     \begin{equation}
%     I_1 = \frac{U_0(s_+(t,x))}{\sqrt{2\pi}} e^{\phi(s_+, t, x)} \cdot \frac{1}{\sqrt{t \lambda K''(s_+)}} \bigg( 
%     1 + \Ocal\bigg(\frac{e^{-\frac{t(\varepsilon^2 \lambda K''(q_0))}{2}} }{\varepsilon \sqrt{\lambda tK''(q_0)}}\bigg) \bigg)\\
%     \end{equation}
    
% $\blacktriangleright$ \ul{The terme $I_2$}

% We have 
%     \begin{equation*}
%         I_2 = \frac{1}{2i \pi} \int_{s_+(t,x) - i\varepsilon}^{s_+(t,x) + i\varepsilon} (U_0(s) -  U_0(s_+(t,x)))  e^{\phi(s,t,x)} \ds 
%     \end{equation*}
    
% Here, we want to go back more or less to the previous case. 

%     \begin{align*}
%          |I_2| 
%          &= 
%          \frac{1}{2i \pi} \sup_{|v| \leq \varepsilon} \big[U_0(s_+ + iv) -  U_0(s_+) \big] \int_{s_+(t,x) - i\varepsilon}^{s_+(t,x) + i\varepsilon} \big| e^{\phi(s,t,x)} \big| \ds \\
%          &=
%          \sup_{|v| \leq \varepsilon} \big[U_0(s_+ + iv) -  U_0(s_+) \big] \frac{1}{\sqrt{2\pi}} e^{\phi(s_+, t, x)} \cdot \frac{1}{\sqrt{t \lambda K''(s_+)}} \bigg( 
%     1 + \Ocal\bigg(\frac{e^{-\frac{t(\varepsilon^2 \lambda K''(q_0))}{2}} }{\varepsilon \sqrt{\lambda tK''(q_0)}}\bigg) \bigg)
%     \end{align*}
    
% It only remains to study $  \sup_{|v| \leq \varepsilon} \big[U_0(s_+ + iv) -  U_0(s_+) \big]$ \\

% By Definition
% \begin{align*}
%     \big[U_0(s_+ + iv) -  U_0(s_+) \big]
%     &= \Big| \int_0^\infty  z^{s_+ - 1}u_0(z) (z^{iv} - 1) \dz \Big| \\
%     &\leq \int_0^{\frac{1}{R}} \Big|  z^{s_+ - 1}u_0(z) (z^{iv} - 1) \Big| \dz
%         + \int_{\frac{1}{R}}^R \Big|  z^{s_+ - 1}u_0(z) (z^{iv} - 1) \Big| \dz \\
%         &\qquad \qquad \qquad \qquad
%         + \int_R^{+\infty} \Big|  z^{s_+ - 1}u_0(z) (z^{iv} - 1) \Big| \dz \\
%     &\leq 2 \int_0^{\frac{1}{R}} \Big|  z^{p_0 + \delta - 1}u_0(z) \Big| \dz
%         +  \int_{\frac{1}{R}}^R \Big|  z^{s_+ - 1}u_0(z) (z^{iv} - 1) \Big| \dz \\
%         &\qquad \qquad \qquad \qquad
%         + 2 \int_R^{+\infty} \Big|  z^{q_0 - \delta - 1}u_0(z) \Big| \dz \\
% \end{align*}

% If $R \to + \infty$, then the first and the third terms vanish. Therefore, we only need to consider the integral $ \int_{\frac{1}{R}}^R \Big|  z^{s_+ - 1}u_0(z) (z^{iv} - 1) \Big| \dz$. We have
%     \begin{align*}
%          \int_{\frac{1}{R}}^R \Big|  z^{s_+ - 1}u_0(z) (z^{iv} - 1) \Big| \dz 
%          &\leq 
%           (\varepsilon \log R)\int_{\frac{1}{R}}^R \Big|  z^{s_+ - 1}u_0(z) \Big| \dz \\
%          &\leq 
%           (\varepsilon \log R)\int_{\frac{1}{R}}^R \Big|  (z^{p_0 + \delta - 1} + z^{q_0 - \delta - 1}) u_0(z) \Big| \dz \\
%     \end{align*}

% \subsection{Interpretation}


% We may rewrite 
% $$x u(t,x)\approx A(t,s_+) e^{t F(s_+(t,x))}$$
% where we define $F$ by
% \begin{equation*}
% F(s)={ K}( s)-1-(s-1) K'\left(s\right).
% \end{equation*}

% The behavior of $xu$ depends on the signe of $F$. 

\section{Introduction}



\section{The Mutation Equation : mutation can be deleterious or beneficial}

\subsection{Location of the population}

Consider the equation
\begin{equation}\left\{
    \begin{array}{ll}
	\frac{\partial}{\partial t}u(t,x) = - \lambda u(t,x) + \lambda \int_{0}^{\infty} k\big( \frac{x}{y}\big)u(t,y) \frac{\dy}{y}
	\eqsepv
	& t \in (0, \infty), x \in \RR_+, \lambda \in (0, \infty) \eqfinv
	\\
	u(0, x) = u_0(x)
	& x \in \RR_+
	\eqfinp
    \end{array}
	\right.
    \label{eq:main}
\end{equation}

We assume that the initial data $u_0(x)$ is a function satisfying the following conditions 
$$ u_0 \geq 0 \eqsepv \int_0^\infty u_0(x)(1+x)\diff x < \infty.$$
It follows that the Mellin transform of $u_0$, $U_0(s)= \int_0^\infty x^{s-1}u(t,x)\diff x$ is analytic on the strip $\Re(\alpha) \in (p_0,q_0)$ where $p_0 \leq 1 < 2 \leq q_0$ \\
For the simplicity, we assume that 
$ p_0=-\infty$ and $q_0 = \infty$\\

By applying the Mellin transform for Equation~\eqref{eq:main}, we have
	\begin{equation}
		\frac{\partial}{\partial t} U(t,s) = - \lambda U(t,s) + \lambda K(s) U(t,s)
		\eqfinp
        \label{eq:explicit_reconstruction}
	\end{equation}
If we have some good hypothesis (see Section~\ref{s:explicit}) on $U_0$ and $K$, then we have an explicit formula for the inverse Mellin transform and it follows that
    \begin{equation*}
        U(t,s) = U_0(s) e^{-\lambda t(1 - K(s) ) }
        \eqsepv
        U_0(s) = \Mcal_{u_0}(s)
        \eqfinp
    \end{equation*}

Here, $K$ is a convex function, such that $K(1) = 1$. We denote $s_* = \argmin_{s\in \RR} K(s)$. 

\subsection{Examples of kernel}

\textit{1.} $k_0(x) = \delta_\alpha$ \\
\textit{2.} $k_0(x) = \frac{1}{b-a}\1{(a,b)}$, $a,b > 0$ \\
Then for all $s \neq 0$, 
    \begin{align*}
        K(s) = \frac{1}{b-a}\int_a^b x^{s-1}\diff x = \frac1s \cdot\frac{b^s - a^s}{b-a}
        \eqfinp
    \end{align*}
\textit{3. The Log--normal kernel.} $k_0(x) = \frac{1}{\sqrt{2\pi}\sigma x} e^{-\frac{(\log(x)-\mu)^2}{2\sigma^2}}$.
By the change of variables $y = \log(x)$, we have
    \begin{align*}
        K(s) 
        &= \int_0^\infty \frac{x^{s-2}}{\sqrt{2\pi}\sigma} e^{-\frac{(\log(x)-\mu)^2}{2\sigma^2}} \diff x \\
        &= \int_{-\infty}^\infty \frac{e^{y(s-1)}}{\sqrt{2\pi}\sigma} e^{-\frac{(y-\mu)^2}{2\sigma^2}} \diff y \\
        &=  \int_{-\infty}^\infty \frac{1}{\sqrt{2\pi}\sigma} e^{\frac{-y^2 - \mu^2 + 2\mu y + 2\sigma^2 y(s-1)}{2\sigma^2}} \diff y \\
        &=  \int_{-\infty}^\infty \frac{1}{\sqrt{2\pi}\sigma} e^{\frac{-y^2 - \mu^2 + 2y[\mu + \sigma^2 (s-1)]}{2\sigma^2}} \diff y \\
        &=  \int_{-\infty}^\infty \frac{1}{\sqrt{2\pi}\sigma} 
        e^{-\frac{(y - [\mu + \sigma^2 (s-1)])^2}{2\sigma^2}}
        e^{\frac{2\mu \sigma^2(s-1)+\sigma^4(s-1)^2}{2\sigma^2}}\diff y \\
        &=  \int_{-\infty}^\infty \frac{1}{\sqrt{2\pi}\sigma} 
        e^{-\frac{(y - [\mu + \sigma^2 (s-1)])^2}{2\sigma^2}}
        e^{\mu (s-1)+\frac{\sigma^2}{2}(s-1)^2}\diff y
        \eqfinp
    \end{align*}
Let denote $  := e^{\mu (s-1)+\frac{\sigma^2}{2}(s-1)^2}$, then 
\begin{align*}
    K(s) 
    &=  C(s,\mu,\sigma) \int_{-\infty}^\infty \frac{1}{\sqrt{2\pi}\sigma} 
        e^{-\frac{(y - [\mu + \sigma^2 (s-1)])^2}{2\sigma^2}}\diff y
        \\
        &=  \frac{C(s,\mu,\sigma)}{\sqrt{2\pi}\sigma} \int_{-\infty}^\infty
        e^{-\frac{(y - [\mu + \sigma^2 (s-1)])^2}{2\sigma^2}}\diff y
\end{align*}



\subsection{A first lemma}

\begin{lemma}Let $k_0$ be a positive measure on $\RR_+$ and $K(s):\RR \to [0,+\infty]$ given by $$K(s)=\int_0^\infty x^{s-1}k_0(\diff x)$$ its Mellin transform. We note $<p_0,q_0>$the fundamental strip of $K$.
\begin{enumerate}
    \item If $K$ is strictly decreasing, then $q_0 = \infty$ and  $\supp k_0 \subset [0,1]$.
    \item If $K$ is strictly increasing, then $p_0 = -\infty$ and $\supp k_0 \subset [1, \infty]$.
\end{enumerate}
\end{lemma}
\begin{proof}
\textit{2st case}: If $K$ is strictly increasing, then $K(s)<K(1)$ for all $s < 1$.


For any $s \in (-\infty, q_0)$, we have
\begin{align*}
    K(s) = \int_0^1 x^{s-1} k_0(\diff x) + \int_1^\infty x^{s-1} k_0(\diff x)
\end{align*}
where these two integrals exist because the integrands are positives and that the integrals are bounded by $K(s)$. \\

Assume there exist a measurable set $A \in \Bcal(\RR_+) \cap (0,1-\varepsilon)$ for some $0 < \varepsilon < 1$ such that $k_0(A) > 0$.  It follows that 
\begin{align*}
   K(1) > K(s) > \int_A x^{s-1} k_0(\diff x)
   > (1-\varepsilon)^{s-1} k_0(A)
   \,,\quad \forall s < 1
   \eqfinp
\end{align*}
But $$(1-\varepsilon)^{s-1} k_0(A) \to \infty \mtext{as} s \to -\infty$$
It's absurd. Therefore for all measurable sets $A \subset (0,1-\varepsilon)$, $k_0(A) = 0$. It follows that  
$$\forall A \in \Bcal\big((0,1)\big), k_0(A) = 0.$$
Therefore, $\supp(k_0) \subset [1, \infty]$.
\end{proof}

The next proposition shows that the behaviour of the repartition of the population depends on $s_* < 1$ or $s_* > 1$. 
\begin{figure}
    \centering
    \includegraphics[scale=0.8]{pictures/severalK.png}
    \caption{The Mellin transform $K(s)$ for two kernel $k_0$. The red point shows the minimum of $K$}
    \label{fig:severalK}
\end{figure}

\begin{proposition} Let $R>0$.
    \begin{enumerate}
        \item Suppose that $s_* < 1$, then
            \begin{align*}
            \lim_{t\to \infty} \int_R^\infty u(t,x) \diff x = \int_0^\infty u_0(x) \diff x
            \eqfinp
            \end{align*}
        \item Suppose that $s_* > 1$, then
            \begin{align*}
            \lim_{t\to \infty} \int_0^R u(t,x) \diff x = \int_0^\infty u_0(x) \diff x
            \eqfinp
            \end{align*}
    \end{enumerate}
\end{proposition}

\begin{proof}\text{}\\

\textbf{First case:} Suppose that $s_* < 1$. Let $c$ be a real $>1$. 
By Mellin inversion, we have that
\begin{align*}
    u(t,x) = \frac{1}{2i\pi}\int_{c-i\infty}^{c+i\infty} U_0(s) x^{-s} e^{-\lambda t(1 - K(s))} \diff s
\end{align*}

It follows that
\begin{align*}
    \int_R^\infty u(t,x) \diff x 
    &= \frac{1}{2i\pi} \int_R^\infty \int_{c-i\infty}^{c+i\infty} U_0(s) x^{-s} e^{-\lambda t(1 - K(s))} \diff s \diff x \\
    &= \frac{1}{2i\pi} \int_{c-i\infty}^{c+i\infty} \int_R^\infty U_0(s) x^{-s} e^{-\lambda t(1 - K(s))} \diff x \diff s \\
    &= \frac{1}{2i\pi} \int_{c-i\infty}^{c+i\infty} U_0(s) \frac{R^{-s+1}}{s-1} e^{-\lambda t(1 - K(s))} \diff s \\
    &= \frac{1}{2i\pi} \int_{c-i\infty}^{c+i\infty}\frac{U_0(s)}{s-1} e^{-\lambda t(1 - K(s))} e^{(-s+1)\ln R} \diff s 
\end{align*}
(The exchange between the two integrals come from the fact that $c > 1$.)

We deform the integration contour towards the left and cross the pole $s = 1$. The Cauchy's residue theorem ensures that
\begin{align}
    \int_R^\infty u(t,x) \diff x &= \mathrm{Res(1)} + \frac{1}{2i\pi} \int_{s_*-i\infty}^{s_*+i\infty} \frac{U_0(s)}{s-1} e^{-\lambda t(1 - K(s))} e^{(-s+1)\ln R} \diff s
    \label{eq:cauchy1}
\end{align}

We define
\begin{align*}
    \Rcal(t) = \frac{1}{2i\pi} \int_{s_*-i\infty}^{s_*+i\infty} \frac{U_0(s)}{s-1} e^{-\lambda t(1 - K(s))} e^{(-s+1)\ln R} \diff s
    \eqfinp
\end{align*}
We have 
\begin{align*}
    |\Rcal(t)| 
    &\leq \frac{R^{1-s_*}}{2\pi} \int_{s_*-i\infty}^{s_*+i\infty} \left|\frac{U_0(s) }{s-1}\right| \cdot \left|e^{-\lambda t(1 - K(s))}\right| \diff s \\
    &= \frac{R^{1-s_*}}{2\pi} \int_{s_*-i\infty}^{s_*+i\infty} \left|\frac{U_0(s) }{s-1}\right| e^{\Re[-\lambda t(1 - K(s))]} \diff s
    \eqfinp
\end{align*}
However, for all $s = s_* + i\nu \in (s_*-i \infty, s_* + i \infty)$
\begin{align*}
    \Re[-\lambda t(1 - K(s))] 
    &= \Re[-\lambda t(K(1) - K(s))] \\
    &= -\lambda t \Re\left( K(1) - \int_0^\infty z^{s-1}k_0(z)\diff z  \right) \\
    &=  -\lambda t\left( K(1) - \int_0^\infty \cos(\nu \log z) z^{s_*-1}k_0(z)\diff z  \right) \\
    &\leq -\lambda t \left( K(1) - \int_0^\infty z^{s_*-1}k_0(z)\diff z  \right) \\
    &= -\lambda t \left( K(1) - K(s_*)  \right)
\end{align*}
By Taylor's expansion in $s_*$, there exists a real $\xi \in (s_*, 1)$ such that 
\begin{align*}
    \Re[-\lambda t(1 - K(s))] = - \lambda t \frac{(s_*-1)^2}{2}  K''(\xi) 
\end{align*}
As $ K''(\xi) > 0$, it follows that 
\begin{align*}
    |\Rcal(t)| 
    &\leq \frac{R^{1-s_*}}{2\pi} \int_{s_*-i\infty}^{s_*+i\infty} \left|\frac{U_0(s) }{s-1}\right| e^{\Re[-\lambda t(1 - K(s))]} \diff s
    \eqfinp \\
    &= \frac{R^{1-s_*}}{2\pi} \int_{s_*-i\infty}^{s_*+i\infty} \left|\frac{U_0(s) }{s-1}\right| e^{- \lambda t \frac{(s_*-1)^2}{2}  K''(\xi) } \diff s \\
    &= e^{- \lambda t \frac{(s_*-1)^2}{2}  K''(\xi)} \times \frac{R^{1-s_*}}{2\pi} \int_{s_*-i\infty}^{s_*+i\infty} \left|\frac{U_0(s) }{s-1}\right| \diff s \\
    &\to_{t \to \infty} 0
    \eqfinp
\end{align*}

Using Equation~\ref{eq:cauchy1}, it follows that 

\begin{align*}
    \int_R^\infty u(t,x) \diff x &= \mathrm{Res(1)} + \Rcal(t) \\
    &= \lim_{s \to 1} (s-1)\frac{U_0(s)}{s-1} e^{-\lambda t(1 - K(s))} e^{(-s+1)\ln R}  + \Rcal(t)\\
    &= U_0(1) + \Rcal(t)\\
    &\to_{t \to \infty} \int_0^\infty u_0(x) \diff x
    \eqfinp
\end{align*}

\textbf{Second case:} Suppose that $s_* > 1$. Let $c$ be a real $<1$. \\
By Mellin inversion, we have that
\begin{align*}
    u(t,x) = \frac{1}{2i\pi}\int_{c-i\infty}^{c+i\infty} U_0(s) x^{-s} e^{-\lambda t(1 - K(s))} \diff s
\end{align*}

It follows that
\begin{align*}
    \int_0^R u(t,x) \diff x 
    &= \frac{1}{2i\pi} \int_0^R \int_{c-i\infty}^{c+i\infty} U_0(s) x^{-s} e^{-\lambda t(1 - K(s))} \diff s \diff x \\
    &= \frac{1}{2i\pi} \int_{c-i\infty}^{c+i\infty} \int_0^R U_0(s) x^{-s} e^{-\lambda t(1 - K(s))} \diff x \diff s \\
    &= \frac{1}{2i\pi} \int_{c-i\infty}^{c+i\infty} U_0(s) \frac{R^{-s+1}}{1-s} e^{-\lambda t(1 - K(s))} \diff s \\
    &= \frac{1}{2i\pi} \int_{c-i\infty}^{c+i\infty} \frac{U_0(s)}{1-s} e^{-\lambda t(1 - K(s))} e^{(1-s)\ln R} \diff s \\
\end{align*}
We deform the integration contour towards the right and cross the pole $s = 1$. The Cauchy's residue theorem ensures that
\begin{align}
    \int_0^R u(t,x) \diff x &= -\mathrm{Res(1)} + \frac{1}{2i\pi} \int_{s_*-i\infty}^{s_*+i\infty} \frac{U_0(s)}{1-s} e^{-\lambda t(1 - K(s))} e^{(-s+1)\ln R} \diff s
    \label{eq:cauchy3}
\end{align}

As previously, we define
\begin{align*}
    \Rcal(t) = \frac{1}{2i\pi} \int_{s_*-i\infty}^{s_*+i\infty} \frac{U_0(s)}{1-s} e^{-\lambda t(1 - K(s))} e^{(-s+1)\ln R} \diff s
    \eqfinp
\end{align*}
and we can show that $\lim_{t\to\infty}\Rcal(t) = 0$. Using Equation~\ref{eq:cauchy3}, it follows that 

\begin{align*}
    \int_0^R u(t,x) \diff x &= -\mathrm{Res(1)} + \Rcal(t) \\
    &= -\lim_{s \to 1} (s-1)\frac{U_0(s)}{1-s} e^{-\lambda t(1 - K(s))} e^{(-s+1)\ln R}  + \Rcal(t)\\
    &= U_0(1) + \Rcal(t)\\
    &\to_{t \to \infty} \int_0^\infty u_0(x) \diff x
    \eqfinp
\end{align*}
\end{proof}

\subsection{Location of the moment}

\begin{proposition} Let $n>1$ be an integer. Suppose that $K$ is minimum in $s_* > 1$, then for all $R > 0$
    \begin{enumerate}
        \item If $s_* > n$ then 
    \begin{align*}
        \lim_{t\to \infty} \int_0^R x^{n-1} u(t,x) \diff x = 0
        \eqfinp
    \end{align*}
        \item If $s_* < n$ then 
        \begin{enumerate}
            \item If $K(n) < 1$
            \begin{align*}
                \lim_{t\to \infty} \int_R^{\infty} x^{n-1} u(t,x) \diff x = 0
                \eqfinp
            \end{align*}
            \item If $K(n) > 1$
            \begin{align*}
                \lim_{t\to \infty} \int_R^{\infty} x^{n-1} u(t,x) \diff x = \infty
                \eqfinp
            \end{align*}
        \end{enumerate}
    \end{enumerate}
\end{proposition}

\begin{proof}
\textbf{First case:}

Suppose that $s_* > n$. Let $c$ be a real $<n$. \\
    By Mellin inversion, we have that
    \begin{align*}
        u(t,x) = \frac{1}{2i\pi}\int_{c-i\infty}^{c+i\infty} U_0(s) x^{-s} e^{-\lambda t(1 - K(s))} \diff s
    \end{align*}

It follows that
\begin{align*}
    \int_0^R x^{n-1}u(t,x) \diff x 
    &= \frac{1}{2i\pi} \int_0^R \int_{c-i\infty}^{c+i\infty} U_0(s) x^{-s+n-1} e^{-\lambda t(1 - K(s))} \diff s \diff x \\
    &= \frac{1}{2i\pi} \int_{c-i\infty}^{c+i\infty} \int_0^R U_0(s) x^{-s+n-1} e^{-\lambda t(1 - K(s))} \diff x \diff s \\
    &= \frac{1}{2i\pi} \int_{c-i\infty}^{c+i\infty} U_0(s) \frac{R^{-s+n}}{n-s} e^{-\lambda t(1 - K(s))} \diff s \\
    &= \frac{1}{2i\pi} \int_{c-i\infty}^{c+i\infty} \frac{U_0(s)}{n-s} e^{-\lambda t(1 - K(s))} e^{(-s+n)\ln R} \diff s
\end{align*}
We deform the integration contour towards the right and cross the pole $s = n$. The Cauchy's residue theorem ensures that
\begin{align}
    \int_0^R x^{n-1}u(t,x) \diff x &= -\mathrm{Res(n)} + \frac{1}{2i\pi} \int_{s_*-i\infty}^{s_*+i\infty} \frac{U_0(s)}{n-s} e^{-\lambda t(1 - K(s))} e^{(-s+n)\ln R} \diff s 
    \label{eq:cauchyn}
\end{align}


As previously, we define
\begin{align*}
    \Rcal(t) = \frac{1}{2i\pi} \int_{s_*-i\infty}^{s_*+i\infty} \frac{U_0(s)}{n-s} e^{-\lambda t(1 - K(s))} e^{(-s+n)\ln R} \diff s
    \eqfinp
\end{align*}
and we can show that $\lim_{t\to\infty}\Rcal(t) = 0$. Using Equation~\ref{eq:cauchyn}, and the fact that $K(n) < 1$ it follows that 


\begin{align*}
    \int_0^R x^{n-1}u(t,x) \diff x &= -\mathrm{Res(n)} + \Rcal(t) \\
    &= -\lim_{s \to n} (s-n)\frac{U_0(s)}{n-s} e^{-\lambda t(1 - K(s))} e^{(-s+n)\ln R}  + \Rcal(t)\\
    &= U_0(n)e^{-\lambda t(1 - K(n))} + \Rcal(t)\\
    &\to_{t \to \infty} 0
    \eqfinp
\end{align*}

\textbf{Second case:}

Suppose that $s_* < n$. Let $c$ be a real $>n$. \\
    By Mellin inversion, we have that
    \begin{align*}
        u(t,x) = \frac{1}{2i\pi}\int_{c-i\infty}^{c+i\infty} U_0(s) x^{s} e^{-\lambda t(1 - K(s))} \diff s
    \end{align*}

It follows that
\begin{align*}
    \int_0^R x^{n-1}u(t,x) \diff x 
    &= \frac{1}{2i\pi} \int_R^{\infty} \int_{c-i\infty}^{c+i\infty} U_0(s) x^{-s+n-1} e^{-\lambda t(1 - K(s))} \diff s \diff x \\
    &= \frac{1}{2i\pi} \int_{c-i\infty}^{c+i\infty} \int_R^{\infty} U_0(s) x^{-s+n-1} e^{-\lambda t(1 - K(s))} \diff x \diff s \\
    &= \frac{1}{2i\pi} \int_{c-i\infty}^{c+i\infty} U_0(s) \frac{R^{-s+n}}{s-n} e^{-\lambda t(1 - K(s))} \diff s \\
    &= \frac{1}{2i\pi} \int_{c-i\infty}^{c+i\infty} \frac{U_0(s)}{s-n} e^{-\lambda t(1 - K(s))} e^{(-s+n)\ln R} \diff s
\end{align*}
We deform the integration contour towards the left and cross the pole $s = n$. The Cauchy's residue theorem ensures that
\begin{align}
    \int_R^{\infty} x^{n-1}u(t,x) \diff x &= \mathrm{Res(n)} + \frac{1}{2i\pi} \int_{s_*-i\infty}^{s_*+i\infty} \frac{U_0(s)}{s-n} e^{-\lambda t(1 - K(s))} e^{(-s+n)\ln R} \diff s 
    \label{eq:cauchyn}
\end{align}


As previously, we define
\begin{align*}
    \Rcal(t) = \frac{1}{2i\pi} \int_{s_*-i\infty}^{s_*+i\infty} \frac{U_0(s)}{s-n} e^{-\lambda t(1 - K(s))} e^{(-s+n)\ln R} \diff s
    \eqfinp
\end{align*}
and we can show that $\lim_{t\to\infty}\Rcal(t) = 0$. Using Equation~\ref{eq:cauchyn}


\begin{align*}
    \int_0^R x^{n-1}u(t,x) \diff x &= \mathrm{Res(n)} + \Rcal(t) \\
    &= \lim_{s \to n} (s-n)\frac{U_0(s)}{s-n} e^{-\lambda t(1 - K(s))} e^{(-s+n)\ln R}  + \Rcal(t)\\
    &= U_0(n)e^{-\lambda t(1 - K(n))} + \Rcal(t)\\
\end{align*}

It follows that 
\begin{equation}
    \int_0^R x^{n-1}u(t,x) \diff x \to_{t \to \infty}
    \begin{cases}
    0 \text{ if $K(n) < 1$}\\
    \infty \text{ if $K(n) > 1$}
    \end{cases}
    \eqfinp
\end{equation}

\end{proof}

\begin{remark}
    We can change the path of integration because : 
\end{remark}
\begin{proof}
    Let $M > 0$. For any $s_* < c$, we have that 
    \begin{align*}
        \bigg| \int_{s_*+iM}^{c+iM} \frac{U_0(s)}{n-s} e^{-\lambda t(1 - K(s))} &e^{(-s+n)\ln R} \diff s \bigg|
        \\
        &\leq 
     \int_{c+iM}^{s_*+iM} \left|\frac{U_0(s)}{n-s}\right| e^{-\lambda t(1 - K(s_*))} e^{\Re((-s+n))\ln R} \diff s \\
     &= 
     \int_{c}^{s_*} \left|\frac{U_0(s+iM)}{n-s-iM}\right| e^{-\lambda t(1 - K(s_*))} e^{(-s+n)\ln R} \diff s.
    \end{align*}
But, for all $s \in (s_*, c)$
\begin{align*}
    \left|\frac{1}{n-s-iM}\right| = \left|\frac{n-s + iM}{(n-s)^2+M^2}\right| \leq \frac{|n-s|}{M^2} + \frac1M
\end{align*}
It follows that $$ \bigg| \int_{s_*+iM}^{c+iM} \frac{U_0(s)}{n-s} e^{-\lambda t(1 - K(s))} e^{(-s+n)\ln R} \diff s \bigg| \to 0$$ when $M \to \infty$.
\end{proof}

\subsection{An explicit solution}
\label{s:explicit}

We assume that     \begin{equation}
    u_0 \geq 0 \mtext{and} \int_0^\infty u_0(x) (1+x) \dx < \infty
    \eqfinp
    \label{H:u0_L1}
\end{equation}

Define 
\begin{equation}
    I(u_0) := \left\{ p \in \RR \eqsepv \int_0^\infty u_0(x) x^{p-1} \dx < \infty \right\}
    \eqfinp
\end{equation}
By Assumption~\eqref{H:u0_L1}, $[1, 2] \subset I(u_0)$. Moreover $I(u_0)$ is a interval of $\RR$.\\
It leads us to define
\begin{subequations}
\begin{align*}
    p_0 &= \inf \{p : p \in I(u_0) \} \in [-\infty, 1] \eqfinv\\
    q_0 &= \sup \{q : q \in I(u_0) \} \in [2, \infty] \eqfinp
\end{align*}    
\end{subequations}

For the simplicity, we assume that 
$ p_0=-\infty$ and $q_0 = \infty$

We also assume that 
\begin{equation}
    \forall \nu \in (p_0, q_0) 
    \eqsepv
    \int_0^\infty |u_0'(x)| x^\nu \dx < \infty
    \eqsepv
    \lim_{x \to 0} x^\nu u_0(x) = 0
    \eqsepv
    \lim_{x \to \infty} x^\nu u_0(x) = 0
    \label{H:DoumicEsco22}
\end{equation}
\begin{equation}
    \forall \nu \in (p_0, q_0) 
    \eqsepv
    \int_0^\infty |u_0''(x)| x^\nu \dx < \infty
    \eqsepv
    \lim_{x \to 0} x^{\nu + 1} u_0'(x) = 0
    \eqsepv
    \lim_{x \to \infty} x^{\nu + 1} u_0'(x) = 0
    \label{H:DoumicEsco23}
\end{equation}

\begin{theorem}(Existence of an explicit solution) Assume \eqref{H:u0_L1}, \eqref{H:DoumicEsco22} and \eqref{H:DoumicEsco23}.
    Then, there exists a unique function $u \in \Ccal([0, \infty); L^1((1+x)dx)) \cap \Ccal^1((0, \infty); L^1((1+x)dx))$ that satisfies Equation~\eqref{eq:main} for all $t > 0$ and almost every $x > 0$. This solution is given by \eqref{eq:explicit_reconstruction} for all $\nu \in (p_0, q_0)$.
    \end{theorem}
    \begin{proof}
    The result is a consequence of the Mellin inversion formula.
    In order to apply it, we have to prove that for any $\nu \in (p_0, q_0)$, 
        \begin{equation*}
            \int\limits_{\nu - i \infty}^{\nu + i \infty} \bigg| U_0(s) e^{\lambda(K(s) - 1) t} x^{-s} \bigg| \ds < \infty
            \eqfinp
        \end{equation*}
    We know that for all $s \in (\nu - i\infty, \nu + i\infty)$, 
        \begin{equation*}
             \bigg| U_0(s) e^{\lambda(K(s) - 1) t} x^{-s} \bigg|  <   \bigg| U_0(s)  e^{\lambda(K(\nu) - 1) t}  x^{-\nu} \bigg|
        \end{equation*}
    
    It only remains to prove that 
        \begin{equation}
                \int\limits_{\nu - i \infty}^{\nu + i \infty} \big| U_0(s)\big| \ds 
                = 
                \int\limits_{-\infty}^{\infty} \big| U_0(\nu + is)\big| \ds
                < \infty
                \label{eq:DouEsc43}
        \end{equation}
    
    This quantity has already been studied under the same assumptions by Doumic and Escobedo \cite{doumic2016time}. Integration by parts lead to 
        \begin{align*}
            U_0(\nu + i v) 
            &= \int_0^\infty u_0(x) x^{\nu-1}x^{i v}\dx \\
            &= \frac{1}{iv + 1} \big[ u_0(x) x^{\nu} x^{i v} \big]_0^\infty - \frac{1}{iv + 1} \int_0^\infty (u_0(x) x^{\nu-1})' x^{i v+1}\dx \\
            \eqfinp
        \end{align*}
    Assumption \eqref{H:DoumicEsco22} guarantees that $[ u_0(x) x^{\nu} x^{i v} ]_0^\infty = 0$. Therefore
        \begin{align*}
            U_0(\nu + i v) 
            &= - \frac{1}{iv + 1} \int_0^\infty (u_0(x) x^{\nu-1})' x^{i v+1}\dx \\
            &= - \frac{\nu - 1}{iv + 1} \int_0^\infty u_0(x) x^{\nu-2} x^{i v + 1}\dx  - \frac{1}{iv + 1} \int_0^\infty u_0'(x) x^{\nu-1} x^{i v + 1}\dx \\
            &= - \frac{\nu - 1}{iv + 1} \int_0^\infty u_0(x) x^{\nu-1} x^{i v}\dx  - \frac{1}{iv + 1} \int_0^\infty u_0'(x) x^{\nu} x^{i v}\dx \\
            &= - \frac{\nu - 1}{iv + 1} U_0(\nu + i v)   - \frac{1}{iv + 1} \int_0^\infty u_0'(x) x^{\nu} x^{i v}\dx
            \eqfinp
        \end{align*}
    It follows that 
        \begin{align*}
            U_0(\nu + i v)
            &= - \frac{1}{(iv + \nu)}\int_0^\infty u_0'(x) x^{\nu} x^{i v}\dx
            \eqfinp
        \end{align*}
    \onlyME{It follow that 
        \begin{equation*}
            |U_0(\nu + i v)| \leq \frac{1}{\sqrt{\nu^2 + v^2}} \int_0^\infty |u_0'(x)| x^{\nu} \dx
            \eqfinp
        \end{equation*}}
    The same calculations and Assumption\eqref{H:DoumicEsco23} allow us to obtain 
        \begin{align*}
            U_0(\nu + i v)
            &=  \frac{1}{(iv + \nu)(iv + \nu + 1)}\int_0^\infty u_0''(x) x^{\nu + 1} x^{i v}\dx
            \eqfinp
        \end{align*}
    It follow that 
        \begin{equation*}
            |U_0(\nu + i v)| \leq \frac{1}{\sqrt{(\nu^2 + v^2)((1+\nu)^2 + v^2)}} \int_0^\infty |u_0''(x)| x^{\nu + 1} \dx
            \eqfinp
        \end{equation*}
    
    We have that $\int_0^\infty |u_0''(x)| x^{\nu + 1} \dx < \infty$ from Assumption\eqref{H:DoumicEsco23}. It allows us to deduce \eqref{eq:DouEsc43} \\
    
    Conversely, if $u_0$ satisfies \eqref{H:u0_L1}, \eqref{H:DoumicEsco22} and \eqref{H:DoumicEsco23}, then 
     $$u \in \Ccal([0, \infty); L^1((1+x)dx)) \cap \Ccal^1((0, \infty); L^1((1+x)dx))$$
    where $u$ is defined by \eqref{eq:explicit_reconstruction}. By construction, it is a solution of the equation \eqref{eq:main}.
    \end{proof}

\section{The behaviour of $u(t,x)$}
We introduce the function: 
    \begin{align*}
        w(t,x) = e^t u(t,x)
        \eqfinp
    \end{align*}
It means that
    \begin{equation*}
        \omega(t,x)
        = \int_{c-i\infty}^{c+i\infty} U_0(s) e^{\lambda K(s)t}x^{-s} \ds
        =  \int_{c-i\infty}^{c+i\infty} U_0(s) e^{\phi(s, t, x)} \ds
        \eqfinp
    \end{equation*}
with $\phi(s,t,x) := \lambda K(s) t - s \ln(x) $.\\

In order to apply the steepest descent method, our first step consist in finding the stationary point of $\phi'(\cdot, t, x) = 0$. 

Define $ s_+(t,x) := (K')^{-1} \Big( \frac{\log x}{\lambda t} \Big)$.


\begin{theorem}
    Suppose that  $u_0$ satisfies the conditions : ...
    Let $\varepsilon(t)$ be any function of $t$ such that $\varepsilon(t) \to 0$ but $t\varepsilon^2(t) \to \infty$ as $t \to \infty$
Then for all $\delta > 0$, there exists a function $\omega_\delta$ such that:
\begin{equation*}
    \begin{split}
    u(t,x) = x^{-s_+(t,x)}e^{(K(-s_+(t,x))-1)t} \times 
    \left( \frac{U_0(s_+(t,x)+\omega_\delta(\varepsilon))}{\sqrt{2\pi tK''(s_+)}} \Theta_1(t)\right)
    \end{split}
\end{equation*}
as $t\to \infty$, uniformly in $x$ such that $-A < s_+(t,x)< A$ and where 
\begin{align*}
    \Theta_1(t) = \left(1 + \Ocal \left( \frac{e^{-\frac{t\varepsilon^2(t)}{2}K''(s_+)}}{\sqrt{t\varepsilon^2(t)K''(s_+)}} \right) \right)
\end{align*}
as $t \to \infty$
\end{theorem}

\begin{proof}
We want to apply the \textit{steepest descent method} \cite{debye1909naherungsformeln} when $t \to \infty$. For this purpose, we split $w(t,x)$ in two parts : 
\begin{align*}
    \omega(t,x)
        &= \int_{s_+-i\infty}^{s_++i\infty} U_0(s)e^{\phi(s, t, x)} \ds \\
        &= \int_{s_+-i\varepsilon}^{s_++i\varepsilon} U_0(s) e^{\phi(s, t, x)} \ds 
        +\int_{\Re(s) = s_+, |Im(s_+)| \geq \varepsilon} U_0(s) e^{\phi(s, t, x)} \ds \\
\end{align*}
We want to show that the "whole mass" comes from the first term. Close to $s_+(t,x)$, we consider the quantity
    \begin{align*}
    \frac{1}{2i \pi} \int_{s_+(t,x) - i\varepsilon}^{s_+(t,x) + i\varepsilon} U_0(s) e^{\phi(s,t,x)} \ds 
    &=
    \frac{1}{2i \pi} U_0(s_+(t,x)) \int_{s_+(t,x) - i\varepsilon}^{s_+(t,x) + i\varepsilon} e^{\phi(s,t,x)} \ds
    \\
    & \qquad \qquad \qquad+ 
    \frac{1}{2i \pi} \int_{s_+(t,x) - i\varepsilon}^{s_+(t,x) + i\varepsilon} (U_0(s) -  U_0(s_+(t,x)))  e^{\phi(s,t,x)} \ds \\
    &= I_1 + I_2
    \end{align*}

\paragraph{\textit{Step 0 -- A first lemma:}}\mbox{}\\
We have
\begin{equation}
    \begin{split}
    K(s) = K(s_+) + (s - s_+(t,x)) K'(s_+(t,x)) + \frac{(s - s_+(t,x))^2}{2} K''(s_+(t,x)) \\+ \frac{(s - s_+(t,x))^3}{6} K'''(c(t,x))
    \end{split}
\end{equation}
where $s < c(t,x) < s_+(t,x)$.
    \begin{align*}
         |K'''(c)|
         &= |\int_0^\infty (\log z)^3 k_0(z) e^{(c-1)\log(z)}\dz|
         \leq \int_0^\infty |(\log z)^3| k_0(z)  e^{\Re (c-1)\log(z)}\dz \\
         &\leq  \boxed{\int_0^\infty |(\log z)^3| k_0(z)  e^{(s_+-1)\log(z)}\dz < \infty}
    \end{align*}
    
Given $\varepsilon > 0$, it follows that for all $t > 0$ and for all $x>0$ such that $|s - s_+| < \varepsilon$, 
    \begin{equation}
        \begin{split}
        \Re K(s) = K(s_+) + \frac{(s - s_+(t,x))^2}{2} K''(s_+(t,x)) + O(\varepsilon^3)
        \eqfinp
        \end{split}
    \end{equation}
and 
    \begin{equation}
        \Re \phi(s, t, x) = \phi(s_+, t, x) + \lambda t \Big(\frac{(s - s_+(t,x))^2}{2} K''(s_+(t,x)) + O(\varepsilon^3)\Big)
        \eqfinp
    \end{equation}
\medskip

\paragraph{\textit{First step -- Study of $I_1$}:}\mbox{}\\
For any $s \in (s_+ - i \infty, s_+ + i \infty)$ such that $|s - s_+| < \varepsilon$.
We know that 
    \begin{equation}
        \Re \phi(s, t, x) = \phi(s_+, t, x) + t \Big(\frac{(s - s_+(t,x))^2}{2} K''(s_+(t,x)) + O(\varepsilon^3)\Big)
        \eqfinp
    \end{equation}

Then 
\begin{align*}
   \int_{s_+(t,x) - i\varepsilon}^{s_+(t,x) + i\varepsilon} e^{\Re \phi(s,t,x)} \ds
    &=
    e^{\phi(s_+, t, x)} \cdot \int_{s_+(t,x) - i\varepsilon}^{s_+(t,x) + i\varepsilon} e^{t \Big(\frac{(s - s_+(t,x))^2}{2} K''(s_+(t,x)) + O(\varepsilon^3)\Big)} \ds \\
    &= 
     i e^{\phi(s_+, t, x)} \cdot \int_{- \varepsilon}^{\varepsilon} e^{t \Big(\frac{(s - s_+(t,x))^2}{2} K''(s_+(t,x)) + O(\varepsilon^3)\Big)} \ds \\
    &= 
     i e^{\phi(s_+, t, x)} \cdot \int_{- \varepsilon}^{\varepsilon} e^{- t \Big(\frac{v^2}{2} K''(s_+(t,x)) + O(\varepsilon^3)\Big)} \ds 
\end{align*}

Now we can see that we are back to the classical method of Laplace. The main idea is to transform the previous integral into a Gauss integral to have
    \begin{equation}
    \int_{- \varepsilon}^{\varepsilon} e^{- t \Big(\frac{v^2}{2} K''(s_+(t,x)) + O(\varepsilon^3)\Big)} \ds 
    \sim_{t \to \infty}
    \sqrt{2\pi} \frac{1}{\sqrt{tK''(s_+)}}
    \end{equation}
    
Indeed, 
\begin{align*}
    \int_{- \varepsilon}^{\varepsilon} e^{- t \Big(\frac{v^2}{2} K''(s_+(t,x)) + O(\varepsilon^3)\Big)} 
    &= 
     \frac{1}{\sqrt{tK''(s_+)}} \int_{- \varepsilon \sqrt{tK''(s_+) + \Ocal(\varepsilon})}^{\varepsilon \sqrt{tK''(s_+) + \Ocal(\varepsilon})} e^{- \frac{y^2}{2}}\dy \\
     \intertext{and}
     \int_{- \varepsilon \sqrt{tK''(s_+) + \Ocal(\varepsilon})}^{\varepsilon \sqrt{tK''(s_+) + \Ocal(\varepsilon})} e^{- \frac{y^2}{2}}\dy &= 
    \sqrt{2\pi} - 2 e^{-\frac{t(\varepsilon^2 K''(s_+) + \Ocal(\varepsilon^3))}{2}} \cdot \Ocal\bigg(\frac{1 }{\varepsilon \sqrt{tK''(s_+) + \Ocal(\varepsilon})}\bigg)
    \eqfinp
\end{align*}
It follows that 
    \begin{align*}
     \frac{1}{2i\pi}   \int_{s_+(t,x) - i\varepsilon}^{s_+(t,x) + i\varepsilon} e^{\Re \phi(s,t,x)} \ds 
     &= 
     \frac{1}{2i\pi} i e^{\phi(s_+, t, x)} \cdot \int_{- \varepsilon}^{\varepsilon} e^{- t \Big(\frac{v^2}{2} K''(s_+(t,x)) + O(\varepsilon^3)\Big)} \ds  \\
     &= 
      \frac{1}{2\pi} e^{\phi(s_+, t, x)} \cdot \frac{1}{\sqrt{tK''(s_+)}} \int_{- \varepsilon \sqrt{tK''(s_+) + \Ocal(\varepsilon})}^{\varepsilon \sqrt{tK''(s_+) + \Ocal(\varepsilon})} e^{- \frac{y^2}{2}}\dy \\
     &=
    \frac{1}{2\pi} e^{\phi(s_+, t, x)} \cdot \frac{1}{\sqrt{tK''(s_+)}} \bigg( 
    \sqrt{2\pi} - 2 e^{-\frac{t(\varepsilon^2 K''(s_+) + \Ocal(\varepsilon^3))}{2}} \cdot \Ocal\bigg(\frac{1 }{\varepsilon \sqrt{tK''(s_+) + \Ocal(\varepsilon})}\bigg) \bigg) \\
    \end{align*}
    
Therefore, 
    \begin{equation}
    I_1 = \frac{U_0(s_+(t,x))}{\sqrt{2\pi}} e^{\phi(s_+, t, x)} \cdot \frac{1}{\sqrt{t \lambda K''(s_+)}} \bigg( 
    1 + \Ocal\bigg(\frac{e^{-\frac{t(\varepsilon^2 \lambda K''(s_+))}{2}} }{\varepsilon \sqrt{\lambda tK''(s_+)}}\bigg) \bigg)\\
    \end{equation}

\medskip
\paragraph{\textit{Second step -- Study of $I_2$}:}\mbox{}\\
We have 
    \begin{equation*}
        I_2 = \frac{1}{2i \pi} \int_{s_+(t,x) - i\varepsilon}^{s_+(t,x) + i\varepsilon} (U_0(s) -  U_0(s_+(t,x)))  e^{\phi(s,t,x)} \ds 
    \end{equation*}
    
Here, we want to go back more or less to the previous integral

    \begin{align*}
         |I_2| 
         &= 
         \frac{1}{2i \pi} \sup_{|v| \leq \varepsilon} \big[U_0(s_+ + iv) -  U_0(s_+) \big] \int_{s_+(t,x) - i\varepsilon}^{s_+(t,x) + i\varepsilon} \big| e^{\phi(s,t,x)} \big| \ds \\
         &=
         \sup_{|v| \leq \varepsilon} \big[U_0(s_+ + iv) -  U_0(s_+) \big] \frac{1}{\sqrt{2\pi}} e^{\phi(s_+, t, x)} \cdot \frac{1}{\sqrt{t \lambda K''(s_+)}} \bigg( 
    1 + \Ocal\bigg(\frac{e^{-\frac{t(\varepsilon^2 \lambda K''(s_+))}{2}} }{\varepsilon \sqrt{\lambda tK''(s_+)}}\bigg) \bigg)
    \end{align*}
    
It only remains to study $  \sup_{|v| \leq \varepsilon} \big[U_0(s_+ + iv) -  U_0(s_+) \big]$ \\

By Definition
\begin{align*}
    \big[U_0(s_+ + iv) -  U_0(s_+) \big]
    &= \Big| \int_0^\infty  z^{s_+ - 1}u_0(z) (z^{iv} - 1) \dz \Big| \\
    &\leq \int_0^{\frac{1}{R}} \Big|  z^{s_+ - 1}u_0(z) (z^{iv} - 1) \Big| \dz
        + \int_{\frac{1}{R}}^R \Big|  z^{s_+ - 1}u_0(z) (z^{iv} - 1) \Big| \dz \\
        &\qquad \qquad \qquad \qquad
        + \int_R^{+\infty} \Big|  z^{s_+ - 1}u_0(z) (z^{iv} - 1) \Big| \dz \\
\end{align*}

If $R \to + \infty$, then the first and the third terms vanish. Therefore, we only need to consider the integral $ \int_{\frac{1}{R}}^R \Big|  z^{s_+ - 1}u_0(z) (z^{iv} - 1) \Big| \dz$. We have
    \begin{align*}
         \int_{\frac{1}{R}}^R \Big|  z^{s_+ - 1}u_0(z) (z^{iv} - 1) \Big| \dz 
         &\leq 
          (\varepsilon \log R)\int_{\frac{1}{R}}^R \Big|  z^{s_+ - 1}u_0(z) \Big| \dz \\
         &\leq 
          (\varepsilon \log R)\int_{0}^{\infty} z^{s_+ - 1}u_0(z)  \dz \\
    \end{align*}

\medskip
\paragraph{\textit{Third step -- A second lemma}:}\mbox{}\\
Outside the neighborhood of $s_+(t,x)$ : $|s - s_+(t,x)| < \varepsilon$. We have
\begin{align*}
    \Re(\phi(s_+ + i \nu, t, x) - \phi(s_+, t, x)) 
    &=
    \Re \bigg( \int_0^\infty x^{s_+ - 1}k_0(x) (x^{i\nu} - 1) \dx \bigg) \\
    &=
    \Re \bigg( \int_0^\infty x^{s_+ - 1}k_0(x) (e^{i\nu \log(x)}- 1) \dx \bigg) \\ 
    &= 
    \int_0^\infty x^{s_+ - 1}k_0(x) [\cos(\nu \log(x))- 1] \dx  \\
    &=: G(v) \leq 0
    \eqfinp
\end{align*}

For any $\nu \in \RR \setminus \{0\}$, we have
\begin{equation}
 \int_0^\infty x^{s_+ - 1}k_0(x) [\cos(\nu \log(x))] \dx < \int_0^\infty x^{s_+ - 1}k_0(x) \dx
 \eqfinp
\end{equation}

It follows that $G(\nu) < 0$, then there exists $\varepsilon > 0$, $R >0$ and $\delta > 0$
such that $\sup_{\varepsilon < \nu < R}G(\nu) < - \delta$.\\

It remains to verify that $\lim_{\nu \to +\infty}G(\nu) \neq 0$. 
\begin{align*}
    G(\nu) 
    &=  \int_0^\infty x^{s_+ - 1}k_0(x) [\cos(\nu \log(x))- 1] \dx \\
    &=  \int_{-\infty}^\infty e^{-s_+ y}k_0(e^{-y})(\cos(-\nu y) - 1 )\dy
    \end{align*}
From Riemann–Lebesgue lemma, we have that 
\begin{equation}
    \lim_{\nu \to +\infty}G(\nu) = - \int_{-\infty}^\infty e^{-s_+ y}k_0(e^{-y})\dy = - \int_{0}^\infty x^{s_+ - 1}k_0(x)\dx < 0
\end{equation}

Therefore
\begin{lemma} For all $\varepsilon > 0$, there exist $\gamma(\varepsilon)$ such
\begin{equation}
    \sup_{|\nu| \geq \varepsilon } \Re(\phi(s_+ + i \nu, t, x)) \leq \Re(\phi(s_+, t, x)) - \gamma(\varepsilon)t
    \eqfinp
\end{equation}
and
\begin{equation}
    \sup_{|\nu| \leq \varepsilon } \Re(\phi(s_+ + i \nu, t, x)) \leq \Re(\phi(s_+, t, x))
    \eqfinp
\end{equation}
\end{lemma}

\medskip
\paragraph{\textit{Fourth step -- Control $\int_{\Re(s) = s_+, |Im(s_+)| \geq \varepsilon} U_0(s) e^{\phi(s, t, x)} \ds$}:}\mbox{}\\
We have that 
\begin{align*}
    \int_{\Re(s) = s_+, |Im(s_+)| \geq \varepsilon} U_0(s) e^{\phi(s, t, x)}
    \ds 
    &\leq  e^{\Re(\phi(s_+, t, x)) - \gamma(\varepsilon)t} \int_{\Re(s) = s_+, |Im(s_+)| \geq \varepsilon} U_0(s) \ds \\
\end{align*}
It follows that 
\begin{align*}
    \int_{\Re(s) = s_+, |Im(s_+)| \geq \varepsilon} U_0(s) e^{\phi(s, t, x)} 
    \ds = \Ocal \left( e^{\Re(\phi(s_+, t, x))} \right)
\end{align*}
\end{proof}


\section{Addition of a term of growth}

Consider the equation
\begin{equation*}\left\{
    \begin{array}{ll}
	\frac{\partial}{\partial t}u(t,x) = - \lambda u(t,x) + \lambda \int_{0}^{\infty} k\big( \frac{x}{y}\big)u(t,y) \frac{\dy}{y} + xu(t,x)
	\eqsepv
	& t \in (0, \infty), x \in \RR_+, \lambda \in (0, \infty) \eqfinv
	\\
	u(0, x) = u_0(x)
	& x \in \RR_+
	\eqfinp
    \end{array}
	\right.
\end{equation*}

By applying the Mellin transform, we have
	\begin{equation*}
		\frac{\partial}{\partial t} U(t,s) = - \lambda U(t,s) + \lambda K(s) U(t,s) + U(t,s+1)
		\eqfinp
	\end{equation*}
It follows that
    \begin{equation}
        \frac{\partial}{\partial t} U(t,s) = - \lambda (1-K(s)) U(t,s) + U(t,s+1)
		\eqfinp
        \label{eq:addgrowth}
    \end{equation}

We consider the Laplace transform of $t \mapsto U(t,s)$
$$ L_s(p) = \int_0^\infty U(t,s)e^{-pt}\diff t$$ 

By applying the Laplace transform to Equation \eqref{eq:addgrowth} we have 
\begin{align*}
    pL_s(p) - U_0(s) = -  \lambda (1-K(s)) L_s(p) + L_{s+1}(p)
\end{align*}
It follows that 
\begin{align*}
    L_{s+1}(p) = - U_0(s) + \big[\lambda (1-K(s)) + p \big]L_s(p)
    \eqfinp
\end{align*}

\begin{remark}
We denote $\alpha(s) = \lambda (1-K(s)) + p$. If $s \in \NN$, and if all quantities are well defined, we have that
\begin{align*}
    L_{s+1}(p) = -U_0(s) -\sum_{n = 1}^{s} \prod_{j = n}^s \alpha(s)U_0(s-1) + L_0(p) \prod_{n = 0}^{s}\alpha(n)
\end{align*}
\end{remark}



\section{An easier equation with growth}

In this section, we suppose that the number of mutations is proportional to the growth rate $x$. This modeling assumption can be justified given that a mutation corresponds to an error in DNA replication. When a cell grows faster, it finds itself replicating its DNA faster. This increases the number of errors, and therefore the number of mutations.\\

We consider the Cauchy problem:
\begin{equation*}\left\{
    \begin{array}{ll}
	\frac{\partial}{\partial t}u(t,x) = - x\mu u(t,x) + \mu \int_{0}^{\infty} y k\big( \frac{x}{y}\big) u(t,y) \frac{\dy}{y} + xu(t,x)
	\eqsepv
	& t \in (0, \infty), x \in \RR_+, \mu \in (0, \infty) \eqfinv
	\\
	u(0, x) = u_0(x)
	& x \in \RR_+
	\eqfinp
    \end{array}
	\right.
\end{equation*}
Here, $u(t,x)$ represents the density of cells with a grotwh rate $x$ at the time $t$. The quantity $\mu$ represents the mutation rate \blue{(NDR: C'est mal dit... car en vrai c'est surtout une sorte de proba de mutation qui ne dépend pas du taux de croissance.)}\\

In order to obtain a solution to this equation, our strategy is to transform it by successively applying a Mellin transform and then a Laplace transform to obtain an equation that is simpler to solve. \\

By applying the Mellin transform in the $x$ variable, we have
	\begin{equation*}
		\frac{\partial}{\partial t} U(t,s) = - \mu U(t,s+1) + \mu K(s) U(t,s+1) + U(t,s+1)
		\eqfinp
	\end{equation*}
It follows that
    \begin{equation}
        \frac{\partial}{\partial t} U(t,s) = \left[- \mu (1-K(s)) + 1\right] U(t,s+1)
		\eqfinp
        \label{eq:addgrowtheasy}
    \end{equation}

We consider the Laplace transform of $t \mapsto U(t,s)$
$$ L(z,s) = \int_0^\infty U(t,s)e^{-zt}\diff t$$ 

By applying the Laplace transform to Equation \eqref{eq:addgrowtheasy} in the $t$ variable we have 
\begin{align*}
    zL_s(z) - U_0(s) = \left[- \mu (1-K(s)) + 1\right] L_{s+1}(z)
\end{align*}
It follows that 
\begin{align}
    \label{eq:wienerhopfeq}
    zL(z,s) = \left[- \mu (1-K(s)) + 1\right] L(s+1,z) + U_0(s)
    \eqfinp
\end{align} 


In the rest of this section, we use the Wiener Hopf method, implemented in \cite{escobedo2007fundamental}, to solve a similar equation. 
\medskip

For this, our goal is to find for any $z \in \CC$ such that $\Re(z) > 0$ a function $L(\cdot,z)$ analytic in the strip $S = \{\xi \in \CC \mid \Re(\xi) \in (s_0, s_0+1)\}$. 
The analycity of $L(\cdot,z)$ strongly depends on the analyticity properties of $K$.\\



\subsection{Reformulating the equation}

The first step is to transform the Equation~\eqref{eq:wienerhopfeq} on a strip into an equation on $\CC \setminus \RR_+$. For this, we introduce the variables:
    \begin{align}
        \zeta &= T(s) = e^{2i\pi(s-s_0)},\\
        g(z,\zeta) &= L(z,s), \\
        \tilde \varphi(\zeta) &= \left[- \mu (1-K(s)) + 1\right] =: 
        \Phi(s)
    \end{align}

We denote as $D = \{ \zeta \in \CC \mid \zeta = re^{i\theta}, r > 0, 0 < \theta < 2\pi \}$. \\

Here, we consider the equation for the fundamental solution, ie : $u_0(x) = \delta_1(x)$ and $U_0(x) = 1$.
\begin{align}
    \label{eq:WH3.5}
    zg(z, x+i0) = \varphi(x)g(z, x-i0) + 1 \, \qquad \text{for all $x\in\RR_+$}
\end{align}
where, for any $x \in \RR_+$, 
\begin{align*}
    g(z, x+i0) &= \lim_{\varepsilon\to 0^+}  g(z, xe^{i\varepsilon}),\\
    g(z, x-i0) &= \lim_{\varepsilon\to 0^+}  g(z, xe^{i(2\pi-\varepsilon)}),\\
    \varphi(x) &= \lim_{\varepsilon\to 0^+}\tilde \varphi(xe^{i\varepsilon}) 
\end{align*}

\begin{hypothesis}\label{hyp:boundK}
    \begin{enumerate}
        \item We assume that $\mu < 1$ and that $K$ is well-defined on $[0, \infty)$.  \blue{Eventually, we may only assume that $K(s) \leq \frac{C}{s^k}$} 
        \item We suppose there exist a vertical strip $<s_0-\varepsilon_b, s_0+1+\varepsilon_b> \in \CC$ such that 
        $$|K(\Re(s))| < \frac{1}{\mu} - 1
        \mtext{for } \varepsilon_b > 0 \eqsepv s \in <s_0-\varepsilon_b, s_0+1+\varepsilon_b>\eqfinp$$ 
    \end{enumerate}
\end{hypothesis}
We note that this last Assumption is natural. Indeed, as soon as we work with a kernel $k_0$ for which we have deleterious and beneficial mutations, we expect $K$ to be a convex function such that $\lim_{s \to \infty} K(s) = \infty$ and $\lim_{s \to -\infty} K(s) = \infty$. When $\mu$ is small enough, it is possible to find a such interval. [\blue{NDR: Détailler ++, rajouter une figure et discuter d'une valeur de $\mu^* $ critique}]

\begin{theorem}
Under Assumption~\ref{hyp:boundK}, the function $g(z, \cdot)$ is analytic on $D$ for any $z \in \CC \setminus \RR_-$ and
\begin{align*}
    g(z,\zeta) 
    &= 
    \frac{\zeta}{2\pi i}\int_0^\infty  \frac{M(z,\lambda-i0)}{\varphi(\lambda)M(z,\zeta)} \frac{\diff \lambda}{\lambda(\lambda-\zeta)} \\
\end{align*}
where $\lambda_0 \in \CC \setminus \RR_+$ and 
$$ M(z, \zeta) := \exp{\frac{1}{2\pi i }\int_0^\infty \ln\left(\frac{z}{\varphi(\lambda)}\right) \left( \frac{1}{\lambda-\zeta} - \frac{1}{\lambda-\lambda_0}\right) \diff \lambda} \eqfinp$$ 
\end{theorem}
\begin{proof}
\blue{For any $s \in \RR_+$, $K(s) > 0 > 1 - \frac1\mu$. Then, 
for all $\zeta \in \RR_+^*$, $\tilde \varphi(\zeta) > 0$.
Since $\tilde \varphi$ is analytic on the half-plane $\RR_+$, it follows that for all $x \in \RR_+$ 
    $$\varphi(x) = \lim_{\varepsilon \to 0^+} \tilde \varphi(xe^{i\varepsilon}) = \tilde \varphi(x) > 0 \eqfinp$$
}

For any complex number $z = re^{i\theta}$ we can define the branch of the complex logarithm
$$\ln(z) := \ln(|z|) + i\theta \eqsepv -2\pi \leq \theta < 0  \eqfinp$$
Let us denote
    \begin{align*}
        H(z,\zeta) := \frac{1}{2\pi i }\int_0^\infty \ln\left(\frac{z}{\varphi(\lambda)}\right) \left( \frac{1}{\lambda-\zeta} - \frac{1}{\lambda-\lambda_0}\right) \diff \lambda
    \end{align*}
where $\lambda_0 \in \CC \setminus \RR_+$ and \blue{$z \in \CC \setminus \RR_-$}\\

There are three things to check to ensure that $H$ is well-defined for $\zeta \in D$. 

\textit{Step 1: Existence when $\lambda \to 0$:}\\
We have that: 
\begin{align*}
    \lim_{\lambda \to 0}\frac{z}{\varphi(\lambda)} = \lim_{\nu\to-\infty}\frac{z}{-\mu(1-K(s_0+i\nu))+1} = \frac{z}{-\mu+1}
    \eqfinp
\end{align*}
\textit{Step 2: Existence when $\lambda \to \infty$:}
\begin{align*}
    \lim_{\lambda \to \infty}\frac{z}{\varphi(\lambda)} = \lim_{\nu\to\infty}\frac{z}{-\mu(1-K(s_0+i\nu))+1} = \frac{z}{-\mu+1}
    \eqfinp
\end{align*}
\textit{Step 3: For all $\lambda \in \RR$, $|\varphi(\lambda)| > \varepsilon$: }\\
For any $s \in s_0 + i\RR$, we have
    \begin{align*}
        |-\mu(1-K(s))+1| = |\mu K(s) - \mu + 1| = \mu \left|K(s) - \left(1-\frac{1}{\mu}\right)\right|
        \eqfinp
    \end{align*}
By Assumption~\ref{hyp:boundK}, for any $s \in s_0 + i\RR$, $|K(s)| \leq K(s_0) < \frac{1}{\mu} - 1$. It follows that
\begin{align*}
    \forall s \in s_0 + i\RR
    \eqsepv
    \left|K(s) - \left(1-\frac{1}{\mu}\right)\right| \geq \frac{1}\mu - 1 - K(s_0) > 0
    \eqfinp
\end{align*}
It follows that $\lambda \mapsto \ln\Big(\frac{z}{\varphi(\lambda)}\Big)$ is bounded on $(0, \infty)$\\

Therefore
$H(z,\zeta)$ is well defined.\\

The formules of Plemej Sojoltski \ndr{ref ?} ensure that for $x > 0$
    \begin{align*}
        H(z,x+i0) &= \frac{1}{2} \ln\left(\frac{z}{\varphi(x)}\right) + \frac{1}{2\pi i} PV\int_0^\infty \ln\left(\frac{z}{\varphi(\lambda)}\right) \left( \frac{1}{\lambda-\zeta} - \frac{1}{\lambda-\lambda_0}\right) \diff \lambda\\
        H(z,x-i0) &= -\frac{1}{2} \ln\left(\frac{z}{\varphi(x)}\right) + \frac{1}{2\pi i} PV\int_0^\infty \ln\left(\frac{z}{\varphi(\lambda)}\right)\left( \frac{1}{\lambda-\zeta} - \frac{1}{\lambda-\lambda_0}\right) \diff \lambda
    \end{align*}
By substracting these expressions and composing by the exponential, we obtain:
\begin{align*}
    \frac{z}{\varphi(x)}
    = \frac{e^{H(z,x+i0)}}{e^{H(z,x-i0)}} 
    := \frac{M(z,x+i0)}{M(z,x-i0)}.
\end{align*}
where 
\begin{equation}
    M(z,\zeta) = e^{H(z,\zeta)}
    \mtext{for}
    \zeta \in D
    \eqsepv z \in \CC \setminus \RR_- \eqfinp
    \label{def:M}
\end{equation}
By using Equation \eqref{eq:WH3.5}, it follows that 
\begin{align*}
    \frac{z}{\varphi(x)} g(z, x+i0) = g(z, x-i0) + \frac{1}{\cte  \varphi(x)} \, \qquad \text{for all $x\in\RR_+^*$}
\end{align*}
and 
\begin{align*}
    \frac{M(z,x+i0)}{M(z,x-i0)} g(z, x+i0) = g(z, x-i0) + \frac{1}{\cte  \varphi(x)} \, \qquad \text{for all $x\in\RR_+^*$}.
\end{align*}
Then 
\begin{align*}
    M(z,x+i0)g(z, x+i0) = M(z,x-i0)g(z, x-i0) + \frac{M(z,x-i0)}{\cte  \varphi(x)} \, \qquad \text{for all $x\in\RR_+^*$}.
\end{align*}

We want to define the function
\begin{align*}
    W(z, \zeta) = \frac{1}{2\pi i}\int_0^\infty  \frac{M(z,\lambda-i0)}{\varphi(\lambda)}\left( \frac{1}{\lambda-\zeta} - \frac{1}{\lambda-\lambda_0}\right) \diff \lambda 
    \eqfinp
\end{align*}
For this purpose, we study the asymptotic behaviour of $M$ when $\lambda \to 0$ and $\lambda \to \infty$. For this, we verify that the function $f(\lambda) = \ln(\frac{z}{\varphi(\lambda)})$ satisfies the hypothesis of the lemma C2 from \cite{escobedo2007fundamental}
\begin{enumerate}
    \item We verify C5: It's straightforward to prove that 
    $$\int_0^\infty \frac{|f(r e^{i\theta)}|}{1+r^2}\diff r < +\infty$$ because due to Assumption~\ref{hyp:boundK}, $f(re^{i\theta})$ is bounded for any $\theta \in (-2\varepsilon_0, 2\varepsilon_0)$ when $\frac{2\varepsilon_0}{2\pi} < \varepsilon_b$. 
    \item We verify C6.the Riemann-Lebesgue lemma ensures that 
    $$ L_1 = \lim_{\lambda \to 0, \lambda \in C(2\varepsilon_0)} f(\lambda) = \ln(\frac{z}{1-\mu}) \eqfinp$$
    $$ L_2 = \lim_{\lambda \to \infty, \lambda \in C(2\varepsilon_0)} f(\lambda) = \ln(\frac{z}{1-\mu}) \eqfinp$$
    \item We verify C7. We have that $\frac{\diff \lambda}{\diff s} = 2i\pi e^{2i\pi(s-s_0)} = 2i\pi \lambda$. Then 
    $\frac{\diff s}{\diff \lambda} = \frac{1}{2i\pi\lambda}$. 
    If follows that 
    $$f'(\lambda) =- \frac{\varphi'(\lambda)}{\varphi(\lambda)} = \frac{-\mu K'(s)}{1 - \mu + \mu K(s)}\frac{\diff s}{\diff \lambda} = \frac{-\mu K'(s)}{1 - \mu + \mu K(s)} \cdot \frac{1}{2i\pi\lambda}. $$
    If $\lambda \to 0$, then $s \to s_0+i\infty$. By the Riemann-Lebesgue lemma, it follows that $$K'(s) \to 0 \mtext{and} K(s) \to 0$$ then  $|\lambda f'(\lambda)| \to 0$ which ensure that $|f'(\lambda)| = o(\frac{1}{\lambda})$. \\. \\
    If $\lambda \to \infty$, then $s \to s_0-i\infty$. By the Riemann-Lebesgue lemma, it follows that $$K'(s) \to 0 \mtext{and} K(s) \to 0$$ then  $|\lambda f'(\lambda)| \to 0$ which ensure that $|f'(\lambda)| = o(\frac{1}{\lambda})$. \\
\end{enumerate}
Therefore, we can apply the lemma C2 from \cite{escobedo2007fundamental} which give the next lemma. 

\begin{lemma}\label{EscARMA:C2}For any $\lambda \in \CC \setminus C(2\varepsilon_0)$, the function $H(z,\zeta))$ is analytic in the domain $D(\varepsilon_0) = \{ `\zeta \in \Scal; \, \zeta = |\zeta|e^{i\theta}, \, \theta \in (-\varepsilon_0, 2\pi+\varepsilon_0)\}$ and 
$$ H(z, \zeta) = -\frac{L}{2\pi i} \ln(\zeta) + o(\ln(|\zeta|)) \mtext{as} \zeta\to 0,\, \zeta \in D(\varepsilon_0)$$ 
$$ H(z, \zeta) = -\frac{L}{2\pi i} \ln(\zeta) + o(\ln(|\zeta|)) \mtext{as} \zeta\to \infty,\, \zeta \in D(\varepsilon_0)$$ 
with $L = \ln(\frac{z}{1-\mu})$. 
\end{lemma}

By Definition, $ M(z, \zeta) = \exp(H(z, \zeta))$. This function is analytic on the domain $D$. Lemma~\ref{EscARMA:C2} ensures that

\begin{lemma}\label{prop:cptM}The function $\zeta \mapsto M(z,\zeta)$ that was defined and analytic for $\zeta \in D$ has an analytic extension to the domain $D(\varepsilon_0)$ for $\varepsilon_0$ small enough. Moreover, for any $\varepsilon > 0$, there exist $\delta \in (0,1)$ such that, for any $z$ in the region
    $$ \Zcal_\varepsilon := \{ z \in \CC :\, \arg z \in (-2\pi-\varepsilon, \varepsilon)\} \eqfinv$$ we have that 
    \begin{equation}
        |M(z, \zeta)| + |\zeta|\left| \frac{\partial}{\partial \zeta}M(z,\zeta)\right|
        \leq 
        C(z) |\zeta|^\delta \mtext{for} |\zeta| < 2 \eqsepv \zeta \in D(\varepsilon_0)
    \end{equation}
    and 
    \begin{equation}
        |M(z,\zeta)| \leq C(z)|\zeta|^{\delta} \mtext{for} |\zeta| > 2 \eqsepv \zeta \in D(\varepsilon_0)
    \end{equation}
\end{lemma}
\begin{proof}We follow the demonstration of Lemma C1 in \cite{escobedo2007fundamental}. From Lemma~\ref{EscARMA:C2} and from the fact that $M(z, \zeta) = \exp(H(z,\zeta))$ it follows that $M$ has an analytic extension on $D(\varepsilon_0)$.
As $\zeta \to 0$, Lemma~\ref{EscARMA:C2} ensures that 
$$ M(z, \zeta) = \exp\left(-\frac{L}{2\pi i} \ln(\zeta) + o(\ln(|\zeta|)) \right)  \mtext{as} \zeta\to 0,\, \zeta \in D(\varepsilon_0) \eqfinp$$
We can deduce by continuity of $M$ that 
$$|M(z, \zeta)| \leq C(z)|\zeta|^\delta
\mtext{for}  \zeta < 1,\, \zeta \in D(\varepsilon_0)$$ with 
 $\delta = -\frac{\arg(z)}{2\pi}$. We note that $\delta \in (0,1)$ because $\arg(z) \in (-2\pi, 0)$.
Since $M(z, \cdot)$ is analytic, the Cauchy estimates ensures that 
\begin{equation*}
    |M(z, \zeta)| + |\zeta|\left| \frac{\partial}{\partial \zeta}M(z,\zeta)\right|
    \leq 
    C(z) |\zeta|^\delta \mtext{for} |\zeta| < 2 \eqsepv \zeta \in D(\varepsilon_0)
\end{equation*}

As $\zeta \to \infty$, 
\begin{align*}
    |M(z, \zeta)| 
    = 
    \exp\left(-\frac{L}{2\pi i} \ln(\zeta) + o(\ln(|\zeta|)) \right)  
\end{align*}
\end{proof}

\begin{lemma}\label{lm:cptW}For any $z \in \Zcal_\varepsilon$, the function $W(z, \cdot)$ defined in (truc) is analytic on $D$ and there is $\delta \in (0,1)$ and a constant $C(z)$ such that
    $$ |W(z, \zeta)| \leq C(z) 
    \mtext{when} |\zeta| < 1 \eqsepv \zeta \in D$$
    and 
    $$ |W(z, \zeta)| \leq C(z)|\zeta|^\delta \mtext{when} |\zeta| > 1 \eqsepv \zeta \in D$$
\end{lemma}
\begin{proof}Lemma~\ref{prop:cptM} and the fact that $\varphi(\lambda)$ is bounded on $(0,\infty)$ ensure that $W(z,\cdot)$ is well defined and analytic on $D$.\\

\textit{(1st case): }For any $\zeta \in D$ such that $|\zeta| < \frac{1}{2}$, 
\begin{align*}
    W(z, \zeta) = \int_0^{2|\zeta|} (\ldots)\diff \lambda + \int_{2|\zeta|}^1 (\ldots) \diff \lambda + \int_1^\infty (\ldots) \diff \lambda := W_1 + W_2 + W_3 \eqfinp
\end{align*}
We have from Lemma~\ref{prop:cptM} that
    $$ |W_3| \leq C(z) \int_1^\infty \lambda^\delta \frac{1+|\lambda_0|}{|\lambda||\lambda - \lambda_0|} < \infty \eqfinp$$
Since $\zeta \in D$, there is $\varepsilon_0 > 0$ such that $\arg \zeta \in (\varepsilon_0, 2\pi - \varepsilon_0)$. If $\lambda \in (0, 2|\zeta|)$, then $|\lambda-\zeta| \geq \frac{\varepsilon_0}{2}|\zeta|$. It follows that 
    $$ W_1 \leq C(z)\int_0^{2|\zeta|} \lambda^\delta \left(\frac{1}{(\lambda-\lambda_0)}+ \frac{2}{\varepsilon_0 |\zeta|}\right) \leq \frac{C(z)}{|\zeta|} \int_0^{2|\zeta|} \lambda^\delta \diff \lambda \leq C(z)|\zeta|^\delta \eqfinp$$
Similarly, we can prove that $|W_2| < \infty$ by using the fact that for all $\lambda \in (2|\zeta|, 1)$ $$|\lambda-\zeta|\geq\frac{|\lambda|}{2}.$$

\textit{(2nd case):} For any $\zeta \in D$ such that $|\zeta| > \frac{1}{2}$, 
\begin{align*}
    W(z, \zeta) = \int_0^{1} (\ldots)\diff \lambda + \int_1^\infty (\ldots) \diff \lambda \eqfinp
\end{align*} 
For all $\lambda \in (0,1)$, \\
For all $\lambda \in (1, \infty)$, 
   $$\left| \int_1^\infty (\ldots) \diff \lambda \right|
   \leq C \int_1^\infty |\lambda|^\delta \frac{|\zeta|} {|\lambda|\cdot|\lambda-\zeta|} \diff \lambda  
   = C \int_1^\infty |\lambda|^\delta \frac{1} {|\lambda|\cdot|\frac{\lambda}{|\zeta|}-1|} \diff \lambda \eqfinp$$
The change of variables $\lambda = x \zeta$ ensures that 
    $$ 
    \left| \int_1^\infty (\ldots) \diff \lambda \right|
    \leq 
    C \int_1^\infty |x|^\delta |\zeta|^\delta \frac{1} {|\zeta||x|\cdot|x-1|} |\zeta|\diff x \eqfinp
    \leq 
    C |\zeta|^\delta \int_1^\infty |x|^\delta \frac{1}{|x|\cdot|x-1|}\diff x \eqfinp
    \leq C |\zeta|^\delta
    $$
\end{proof}


The Sokhotski–Plemelj formula ensures that 
\begin{align*}
    \frac{M(z,x-i0)}{\varphi(x)}
    = 
    W(z,x+i0) - W(z,x-i0)
    \mtext{for any}
    x > 0
    \eqfinp
\end{align*}
Therefore, 
\begin{align*}
    M(z,x+i0)g(z, x+i0) -W(z,x+i0) = M(z,x-i0)g(z, x-i0) - W(z,x-i0) \, \qquad \text{for all $x\in\RR_+$}.
\end{align*}
It follows that the function $C(z,\cdot)$ defined by 
$$ C(z,\cdot) := M(z,\cdot)g(z,\cdot) - W(z,\cdot)$$
is analytic in $\CC \setminus \{0\}$.\\

Using the fact that $g$ is bounded on $\CC$, Lemma~\ref{prop:cptM} and Lemma~\ref{lm:cptW}
, the Liouville's theorem ensures that
\begin{align*}
    \forall z \in \CC \setminus \RR_-
    \eqsepv
    C(z,\zeta) = C(z)
\end{align*}
Then 
\begin{align}
    g(z,\zeta) = \frac{C(z)+W(z,\zeta)}{M(z,\zeta)}
\end{align}
However, Lemma~\ref{prop:cptM} and the boundeness of $g$ ensure that
\begin{align}
C(z) 
= \lim_{\zeta \to 0}[M(z,\cdot)g(z,\cdot) - W(z,\cdot)] 
= -\lim_{\zeta \to 0}W(z,\zeta)
\end{align}
As, 
$$ \lim_{\zeta \to 0}W(z,\zeta) 
= \frac{1}{2\pi i}\int_0^\infty  \frac{M(z,\lambda-i0)}{\varphi(\lambda)}\left( \frac{1}{\lambda} - \frac{1}{\lambda-\lambda_0}\right) \diff \lambda $$
we have that
\begin{align*}
    C(z) + W(z,\zeta)
    &= 
    \frac{1}{2\pi i}\int_0^\infty  \frac{M(z,\lambda-i0)}{\varphi(\lambda)}\left( \frac{1}{\lambda-\zeta}-\frac{1}{\lambda}\right) \diff \lambda \\
    &= 
    \frac{\zeta}{2\pi i}\int_0^\infty  \frac{M(z,\lambda-i0)}{\varphi(\lambda)} \frac{1}{\lambda(\lambda-\zeta)} \diff \lambda \\
\end{align*}
It follows that 
\begin{align*}
    g(z,\zeta) 
    &= 
    \frac{\zeta}{2\pi i}\int_0^\infty  \frac{M(z,\lambda-i0)}{\varphi(\lambda)M(z,\zeta)} \frac{\diff \lambda}{\lambda(\lambda-\zeta)} \\
\end{align*}
For any $z \in \CC \setminus \RR_-$, Equation~\eqref{def:M} ensures that $M(z, \cdot)$ does not vanish on $D$ and Lemma~\ref{prop:cptM} ensures that $M(z, \cdot)$ is analytic on $D(\varepsilon_0$). It follows that $g(z, \cdot)$ is analytic on $D$ for any $z \in \CC \setminus \RR_-$.

\end{proof}

\subsection{The solution of the Equation}
Now we want to return to the solution of Equation~\eqref{eq:wienerhopfeq}. For this purpose, we revert the change of variables. We have that 
\begin{align}
    L(z,s) = g(z, T(s)) 
    \eqfinp
\end{align}
We define $m(z,s) := M(z, T(s))$. Then

\begin{align*}
    L(z,s) &= g(z,T(s)) \\ 
    &= 
    \frac{T(s)}{2\pi i}\int_0^\infty  \frac{M(z,\lambda-i0)}{\varphi(\lambda)m(z,s)} \frac{\diff \lambda}{\lambda(\lambda-T(s))} \\
    &= 
    \frac{1}{2\pi i}\int_0^\infty  \frac{M(z,\lambda-i0)}{\varphi(\lambda)m(z,s)} \frac{\diff \lambda}{\lambda\Big(\frac{\lambda}{T(s)}-1\Big)} \\
    &= 
    \frac{1}{2\pi i}\int_0^\infty  \frac{M(z,\lambda-i0)}{\varphi(\lambda)m(z,s)} \frac{\diff \lambda}{\lambda\Big(\frac{\lambda}{e^{2i\pi(s-s_0)}}-1\Big)} \\
\end{align*}
The change of variable $\lambda = e^{2i\pi(y-s_0)}$ gives
\begin{align*}
    L(z,s)
    &= 
    - \frac{2\pi i}{2\pi i}\int_{s_0-i\infty}^{s_0+i\infty}  \frac{m(z,y-i0)}{\varphi(e^{2i\pi(y-s_0)})m(z,s)} \frac{\diff y}{(e^{2i\pi(y-s)}-1)} \\
    &= 
   - \int_{s_0-i\infty}^{s_0+i\infty} \frac{1}{\varphi(e^{2i\pi(y-s_0)})}  \frac{m(z,y-i0)}{m(z,s)} \frac{\diff y}{(e^{2i\pi(y-s)}-1)} \\
\end{align*}
where 
$$m(z,s) = M(z, T(s))$$
The change of variable shows that
\begin{equation}
    m(z,s) = \exp \left[
    - \int_{s_0-i\infty}^{s_0+i\infty}\log \left( \frac{z}{\Phi(
    \sigma)}\right)
    \cdot \left(
        \frac{1}{1-e^{2i\pi(s-\sigma)}} - \frac{1}{1-e^{2i\pi(\sigma_0-\sigma)}}
     \right) \diff \sigma
    \right]
\end{equation} 
with $\lambda_0 = e^{2i\pi(y_0 - s_0)}$. 
We have that 
\begin{equation}
    L(z,s) 
    = 
   - \int_{s_0-i\infty}^{s_0+i\infty} \frac{1}{\Phi(y)}  \frac{m(z,{\MD y-i0})}{m(z,s)} \frac{\diff y}{(e^{2i\pi(y-s)}-1)}
\end{equation}

We note that the term $\frac{m(z,y-i0)}{m(z,s)}$ depends on $z$. For reasons that will become clearer later, we want to isolate this dependency in a term that can be calculated explicitly. 
It can be shown that 
\begin{equation*}
    m(z,y) = V(y)
    \exp \left[
    - \int\limits_{s_0-i\infty}^{s_0+i\infty} \left[\log \left( \frac{z}{1-\mu}\right) + 2i\pi n\right]
    \cdot \left(
        \frac{1}{1-e^{2i\pi(y-\sigma)}} - \frac{1}{1-e^{2i\pi(\sigma_0-\sigma)}}
     \right) \diff \sigma
    \right]
\end{equation*}
where 
$$ V(y) := \exp \left[
    - \int\limits_{s_0-i\infty}^{s_0+i\infty}\log \left( \frac{1-\mu}{\Phi(\sigma)}\right)
    \cdot \left(
        \frac{1}{1-e^{2i\pi(y-\sigma)}} - \frac{1}{1-e^{2i\pi(\sigma_0-\sigma)}}
     \right) \diff \sigma 
    \right]  $$ 
and $n \in \NN$.\\

Then 
\begin{equation*}
    \frac{m(z,y)}{m(z,s)} 
    = 
    \frac{V(y)}{V(s)} \times E(z,y,s)
\end{equation*}
where $$E(z,y,s) := \exp \left[
    - \int_{s_0-i\infty}^{s_0+i\infty} \left[\log \left( \frac{z}{1-\mu}\right) + 2i\pi n\right]
    \cdot \left(
        \frac{1}{1-e^{2i\pi(y-\sigma)}} - \frac{1}{1-e^{2i\pi(s-\sigma)}}
     \right) \diff \sigma \right]
     $$

We have that
\begin{equation*}
    E(z,y,s) 
    = 
    \exp
        \left[
     \left(\log \left( \frac{z}{1-\mu}\right) + 2i\pi n\right)
    \cdot (y-s)\right]
\end{equation*}

Then 
\begin{equation}\label{eq:exprL}
    L(z,s) 
    = 
   - \int_{s_0-i\infty}^{s_0+i\infty} \frac{1}{\Phi(y)} \frac{V(y)}{V(s)} e^{ \alpha(z)(y-s)}
   \frac{\diff y}{(e^{2i\pi(y-s)}-1)}
\end{equation}
where $\alpha(z) :=  \left[
    \left(\log \left( \frac{z}{1-\mu}\right) + 2i\pi n\right)
    \right]$
\subsection{Inversion of the transforms}

We aim to apply the inverse Laplace transform and the inverse Mellin transform to $L(z,s)$ to get an expression of $u(t,x)$.\\

We define the inverse Mellin transform of $L(z,\cdot)$ by 
\begin{equation*}
    \Mcal^{-1}(L(z))(x) = \frac{1}{2i\pi} \int_{s_0-i\infty}^{s_0+i\infty} L(z,s)x^{-s}\diff s
    \eqsepv
    z \in \CC \setminus \RR_-
\end{equation*}

We have that 
\begin{align*}
    V(y+1) 
    &= 
    \exp \left[
        - \int_{s_0-i\infty}^{s_0+i\infty}\log \left( \frac{1-\mu}{\Phi(\sigma)}\right)
        \cdot \left(
            \frac{1}{1-e^{2i\pi(y+1-\sigma)}} - \frac{1}{1-e^{2i\pi(\sigma_0-\sigma)}}
         \right) \diff \sigma \right]  \\
\end{align*}
Here, singularities can only come from $(s-\sigma) \in \ZZ$. 

 Let $s \in \CC$ such that

\ndr{A rédiger proprement; De toute évidence, cela devrait faire apparaitre $s_*$ et $s^*$ qui sont des valeurs critiques de l'hypothèse (qu'il faudra définir au dessus)}
 We have the formula 
 \begin{equation}
     V(s+1) = V(s)\frac{\Phi(s+1)}{1-\mu}
     \mtext{for}
     s \in (s_0, s_0+1) + i\RR
     \eqfinp
 \end{equation}
 We note that we have some pope for $s+1 = s_*$.
 We can extend this formula 
 



\bibliographystyle{plain}
\bibliography{oneforall}

\end{document}
